\documentclass[12pt]{book}
\usepackage[utf8]{inputenc}

\usepackage{textcomp}

\usepackage[T1]{fontenc}

\usepackage[utf8]{inputenc}
\usepackage{graphicx}
\usepackage{amsmath}
\usepackage{amsfonts}
\usepackage{amsthm} % unnumbered environments

\usepackage{tikz}
\usetikzlibrary{babel}
\usepackage[spanish]{babel}


\usepackage{lmodern}

\usepackage{minted} % código haskell
\usepackage{tikz-cd} % diagramas conmutativos

\usepackage[hidelinks]{hyperref} % para hiperreferencias

\usepackage{color} %para el color

% pies de página y esas cosas
%\usepackage{fancyhdr}
%\pagestyle{fancy}

% cambiar chapter por capítulo
%\renewcommand{\chaptername}{Capítulo}

% algunas macros que usaré frecuentemente
% categorías genéricas
\newcommand{\C}[0]{\mathcal{C}}
\newcommand{\D}[0]{\mathcal{D}}
\newcommand{\E}[0]{\mathcal{E}}
\newcommand{\Obj}[1]{\mathcal{O}b(#1)}
\newcommand{\Arr}[1]{\mathcal{A}r(#1)}
\newcommand{\arr}[2]{#1\longrightarrow #2}
\newcommand{\nat}[2]{#1\Rightarrow #2}

\DeclareMathOperator{\Hom}{Hom}
\DeclareMathOperator{\Nat}{Nat}

% macros para nombres de categorías
\newcommand{\Set}[0]{\texttt{Set}}
\newcommand{\Grp}[0]{\texttt{Grp}}
\newcommand{\Ring}[0]{\texttt{Ring}}
\newcommand{\Top}[0]{\texttt{Top}}
\newcommand{\Hask}[0]{\texttt{Hask}}
\newcommand{\VectK}[0]{\texttt{Vect-}K}
\newcommand{\Ab}[0]{\texttt{Ab}}

% algebra
\newcommand{\Z}[0]{\mathbb{Z}}
\newcommand{\Q}[0]{\mathbb{Q}}

% short code snippets
\newcommand{\cod}[1]{\textcolor{violet}{\texttt{#1}}}
%\newcommand{\cod}[1]{\mintinline{haskell}{#1}}
\newcommand{\uno}{*}
\newcommand{\final}{*}
\newcommand{\inicial}{\oslash}
\newcommand{\nulo}{\textbf{0}}

\newcommand{\funccat}[2]{#2^{#1}}

\newcommand{\adjoint}{\dashv}

\newcommand{\pcomp}{$\sim$}


% bibliografía
\usepackage{biblatex}
\addbibresource{references.bib}


\setcounter{secnumdepth}{3} % que numere las subsecciones también

\graphicspath{ {images/} }

% ambientes: teoremas, definiciones
\newtheorem{definition}{Definición}
\newtheorem*{definition*}{Definición}
\newtheorem{proposition}{Proposición}
\newtheorem*{proposition*}{Proposición}
\newtheorem{theorem}{Teorema}


\title{
	{Patrones de diseño inspirados por teoría de categorías}\\
	{\large Universidad de Granada}\\
}
\author{Braulio Valdivielso Martínez}
\date{\today}


\usepackage{eso-pic}
\usepackage{graphicx}
%\usepackage{fourier}
%\usepackage[beramono,eulerchapternumbers,linedheaders,parts,a5paper,dottedtoc, manychapters,pdfspacing]{classicthesis}
%\usepackage[beramono, eulerchapternumbers,linedheaders,dottedtoc, manychapters,pdfspacing]{classicthesis}
\usepackage[top=1in, bottom=1.5in, left=1in, right=1in]{geometry}

\newcommand\BackgroundPic{%
	\put(0,0){%
		\parbox[b][\paperheight]{\paperwidth}{%
			\vfill
			\centering
			\includegraphics[width=\paperwidth,height=\paperheight,%
			keepaspectratio]{ugrA4.pdf}%
			\vfill
}}}


\begin{document}

\begin{titlepage}
\AddToShipoutPicture*{\BackgroundPic}

\newlength{\centeroffset}
\setlength{\centeroffset}{-0.5\oddsidemargin}
\addtolength{\centeroffset}{0.5\evensidemargin}
\thispagestyle{empty}

\vspace*{4cm}

\noindent\hspace*{\centeroffset}\begin{minipage}{\textwidth}

	\centering
	%\includegraphics[width=0.4\textwidth]{logougrnuevo.jpg}\\[1.4cm]

	\textsc{ \huge Trabajo Fin de Grado\\[0.2cm]}
	\medskip
	\textsc{ \Large Grado en Matemáticas}\\[0.8cm]
	% Upper part of the page
	%
	% Title

	{\Huge\bfseries \makeatletter \@title \makeatother
	}

\end{minipage}

\vspace{2cm}
\noindent\hspace*{\centeroffset}\begin{minipage}{\textwidth}
	\centering
	\textbf{\large Autor}\\ {\large \makeatletter \@author \makeatother}\\[5.5ex]
	%		\textbf{\large Tutor}\\
	%		{\large Pedro A. García-Sánchez}\\[1.6cm]
	%		\textsc{---}\\
	%		Granada, febrero de 2018
\end{minipage}
%\addtolength{\textwidth}{\centeroffset}
%\vspace{\stretch{2}}

\vfill
\begin{minipage}[t][1cm][t]{1.7cm}
	{\large  \textbf{Tutor:}}
\end{minipage}\hfill
\begin{minipage}[t][1cm][t]{7cm}
	{\large María Burgos Navarro\\
		\emph{Didáctica de la Matemática}}\\[.7em]
	{\large Pedro A. García Sánchez\\
		\emph{Álgebra}}
\end{minipage}\hfill
\begin{minipage}[t][1cm][t]{3.5cm}
	\large{Junio 2018}
\end{minipage}

\end{titlepage}

\tableofcontents

\chapter*{Resumen y palabras clave}


\chapter*{Abstract and keywords}
To mum and dad

\chapter*{Declaration}
I declare that..

\chapter*{Introducción}
El lenguaje de programación haskell surge en el año 1990 con
la idea de aunar los esfuerzos de múltiples instituciones académicas
en el desarrollo de un lenguaje de programación funcional. Este
lenguaje de programación estaba pensado para ser una base
común sobre la que realizar docencia e investigación en la
programación funcional.

Haskell cumplió su propósito de ser el vehículo de numerosa innovaciones
en materia del diseño de lenguajes de programación. Entre otros
desarrollos, haskell hizo de un concepto nacido en la teoría de
categorías, las mónadas, una construcción fundamental para el lenguaje.

Conforme se refinaban las especificaciones del lenguaje, la idea de
utilizar haskell en la industria era cada vez
más factible. Hoy en día haskell
es usado en empresas de sectores tan diversos como la banca,
las redes sociales o la mensajería instantánea.

El objetivo del trabajo será comprender la forma en la que
la teoría de categorías ha sido aplicada al diseño de software
en el lenguaje de programación haskell. El objetivo del trabajo,
naturalmente, no puede cumplirse sin estudiar la teoría de categorías
de la que emanan estas aplicaciones.

Comenzamos el trabajo por las definiciones más básicas de la
teoría de categorías: las categorías y los funtores. Ya en el
primer capítulo podremos encontrar los primeros ejemplos de
conceptos categóricos que se usan
diariamente en el desarrollo de software en haskell.

En el segundo capítulo mostraremos la potencia expresiva de
las categorías para la realización de distintas construcciones
sobre estructuras matemáticas diferentes. Especializaremos algunas
de estas construcciones a la programación.

El tercer capítulo tratará una de las nociones más importantes
de la teoría de categorías, que nos servirá
más adelante para enunciar la definición de mónada:
las transformaciones naturales.
Aprovecharemos la teoría construida hasta el momento
para enunciar y demostrar el lema de Yoneda, uno de los resultados
centrales de la teoría de categorías. En la parte de aplicaciones
a la programación veremos cómo las transformaciones naturales
permiten establecer relaciones entre constructores de tipos
polimórficos.

En el cuarto capítulo conseguiremos definir
qué son las monadas. Tendremos que pasar
antes por la definición de adjunción, una abstracción que ha sido
identificada por múltiples autores como uno de los conceptos clave
de la teoría. Finalizaremos la parte del trabajo dedicada a las
matemáticas analizando la relación entre las adjunciones y las
mónadas. Al final del capítulo intentaremos justificar el rol
de las mónadas en la programación funcional.

Finalizaremos el trabajo en el capítulo quinto. Dedicaremos
este capítulo a la exposición de dos patrones de diseño que,
junto con las mónadas, muestran la influencia de la teoría de
categorías en la forma en la que se programa en haskell. Aprovecharemos
también para comentar brevemente otros patrones de diseño
que indicarían líneas de futuro trabajo.

El libro \textit{Topos, Triples and Theories} ha sido la referencia
principal para el estudio de la teoría de categorías durante
la elaboración de este trabajo. Ocasionalmente
se ha utilizado el libro
\textit{Categories for the working mathematician}
para la consulta de algún tema concreto.

El principal recurso de consulta sobre el lenguaje de programación
haskell ha sido \textit{Real World Haskell}. Esporádicamente se
han consultado otras fuentes que serán referenciadas cuando corresponda
a lo largo del texto.

\tableofcontents



\chapter{Categorías y funtores}
\section{Categorías}
\subsection{Definición}

Una categoría $\mathcal{C}$ queda determinada por dos colecciones:

\begin{enumerate}
\item $\Obj{\C}$ a cuyos elementos nos referiremos como
  \emph{objetos} de $\mathcal{C}$.
\item $\Arr{\C}$ a la que nos referiremos como las
  \emph{flechas} de $\mathcal{C}$. A cada flecha se le puede asignar
  un par de objetos: al primero de ellos le llamaremos \emph{origen}
  (o dominio) y al otro \emph{destino} (o \emph{codominio}). Con la notación
  $f : \arr{A}{B}$ estaremos afirmando que $f$ es una flecha
  que tiene como origen el objeto $A$ y como destino el objeto $B$.
  No supondremos en general que $\mathcal{A}r(\mathcal{C})$ es un
  conjunto pero a lo largo de este texto sí que asumiremos que
  $Hom(A, B) = \{ f: f : A \longrightarrow B \}$ lo es. A las categorías
  en las que $Hom(A, B)$ son conjuntos se les suele llamar categorías
  localmente pequeñas. En este texto solo trataremos categorías
  localmente pequeñas.
\end{enumerate}
Para poder hablar de categorías necesitamos también una operación de composición $\circ$ que funcione de la siguiente
forma:

\begin{enumerate}
\item Para cada terna de objetos $A, B, C$ de la categoría $\mathcal{C}$
  diremos que los pares de flechas de $Hom(B, C)\times Hom(A, B)$ se
  pueden componer y tendremos que
  $\circ : Hom(B, C)\times Hom(A, B) \rightarrow Hom(A, C)$. Dicho
  de otra forma si $f: A\longrightarrow B$ y $g : B \longrightarrow C$
  tenemos que $g \circ f : A \longrightarrow C$.
\item $\circ$ es asociativa en el siguiente sentido: si
  $f : A \longrightarrow B, g: B \longrightarrow C,$
  $h : C \longrightarrow D$ tenemos que
  $(h \circ g) \circ f = h \circ (g \circ f)$.
\item Para cada objecto $C$ de la categoría $\mathcal{C}$ existe una
  flecha a la que llamaremos $1_C : C \longrightarrow C$ tal que
  para cualquier flecha $f: X \longrightarrow C$ tenemos que
  $1_C \circ f = f$ y para cualquier flecha $g : C \longrightarrow Y$
  se cumple $g \circ 1_C = g$ para cualquiera que sean los objetos
  $X, Y$ de $\mathcal{C}$.
\end{enumerate}

\textcolor{red}{Voy a reescribirte esto, creo que la forma en que se introduce una categoría da una idea de la escuela que se sigue o de la filosofía con que se afronta esta teoría. Te voy a introducir en categorías como lo haría yo. TU ELIGES LA FORMA DEFINITIVA. No haré esto con el resto del trabajo pero el comienzo es importante.}

Tradicionalmente las matemáticas están fundamentadas en  \emph{una teoría de conjuntos} y bajo esta fundamentación el concepto de \emph{conjunto} es básico. Entendemos lo que es un conjunto pero no tratamos de dar una definición formal de este. Ocurre lo mismo con los conceptos de \emph{elemento} y \emph{pertenece} que son básicos en esta teoría.  En la actualidad se puede utilizar la \emph{teoría de categorías} para fundamentar las matemáticas y en este sentido los conceptos de \emph{categoría, objeto, flecha} y \emph{composición} serían los conceptos básicos que se intentan entender sin dar una definición formal de estos. Siguiendo esta idea podemos decir que tendremos una categoría $\mathcal{C}$ si:
\begin{itemize}
	\item Conocemos sus objetos, que denotamos $A,B,C,\ldots.$
	\item Conocemos sus flechas, que denotamos $f,g,h,\ldots.$
	\item Para cada flecha $f$ conocemos su dominio $A$ y su codominio $B$, que serán objetos de $\mathcal{C}$, escribiremos $f:A\rightarrow B$ o bien $A\xrightarrow{f} B$.
	\item Para cada dos \emph{flechas componibles} $A\xrightarrow{f} B\xrightarrow{g} C$ conocemos su composición $g\circ f :A\rightarrow C$.
\end{itemize}
Todos estos datos, que determinan una categoría $\mathcal{C}$, tienen que cumplir las siguientes propiedades o axiomas:
\begin{enumerate}
	\item La composición es asociativa, en el siguiente sentido: si $f : A \longrightarrow B, g: B \longrightarrow C$ y $h : C \longrightarrow D$ se ha de cumplir
	$(h \circ g) \circ f = h \circ (g \circ f)$.
	\item Existen identidades, esto es: para cada objecto $C$ existe una
	flecha, a la que llamaremos identidad en $C$ y que denotaremos $1_C : C \longrightarrow C$, que cumple que
	para cualquiera flechas $f: X \longrightarrow C$ y  $g : C \longrightarrow Y$ se tiene
	$1_C \circ f = f$  y  $g \circ 1_C = g$.
\end{enumerate}

Con esta aproximación a la teoría de categorías no haríamos uso de la teoría de conjuntos. Sin embargo vamos a optar por una aproximación \emph{no tan categórica}. En toda esta memoria vamos a asumir que para cada par de objetos $A$ y $B$, de la categoría $\mathcal{C}$, las flechas de $\mathcal{C}$ con dominio $A$ y codominio $B$ forman un conjunto, que denotaremos $Hom_{\mathcal{C}}(A,B)$ o simplemente $Hom(A,B)$. De manera que la composición en $\mathcal{C}$ determina aplicaciones $\circ : Hom(B, C)\times Hom(A, B) \rightarrow Hom(A, C)$ para cada terna de objetos $A,B,C$ de $\mathcal{C}$. En este sentido, en esta memoria trataremos sólo con categorías \emph{localmente pequeñas}. Denotaremos $\Obj{\C}$ y $\Arr{\C}$ a las \emph{clases} de todos los objetos y todas las flechas respectivamente de $\C$.




Mostramos a continuación algunos ejemplos de categorías.

\subsection{Ejemplos}
\paragraph{Conjuntos}
Uno de los más típicos ejemplos de categorías es $\Set$,
la categoría de los conjuntos. En esta categoría
cada conjunto es un objeto (no se puede asumir que
$\mathcal{O}b(\mathcal{C})$ es un conjunto por ejemplos como este:
no tiene sentido hablar del conjunto de todos los conjuntos) y cada
aplicación $f$ entre conjuntos con dominio el conjunto $A$ y como codominio
el conjunto $B$ es una flecha $f : A \longrightarrow B$. La composición
es la composición habitual de aplicaciones y las identidades
$1_C : C \longrightarrow C$ son las aplicaciones identidad en
cada conjunto $C$.
Comprobar que se cumplen los axiomas de las categorías es una tarea
rutinaria.

\paragraph{Otras estructuras matemáticas}
Gran parte de las estructuras que se estudian en matemáticas forman
categorías si consideramos sus morfismos como flechas.
Podemos dar multitud de ejemplos de este tipo:

\begin{itemize}
\item \texttt{Grp}: la categoría en la que los objetos son grupos
  y las flechas son los homomorfismos de grupos.
\item \texttt{Top}: la categoría en la que los objetos son espacios
  topológicos y las flechas son funciones continuas.
\item \texttt{Ring}: la categoría en la que los objetos son
  anillos y las flechas son homomorfismos de anillos.
\end{itemize}
La lista sigue y sigue.

\paragraph{Monoides}
Proponemos este ejemplo para evitar la asumción de que en una categoría
los objetos deben ser estructuras matemáticas y las flechas entre ellos
aplicaciones que preservan la estructura. Definimos una categoría
con un solo objeto al que llamaremos $*$. El conjunto de flechas
será $Hom(*, *) = \mathbb{Z}$ y
$\circ : Hom(*, *)\times Hom(*, *)\rightarrow Hom(*, *)$ quedará definido
por $f \circ g = f + g$ donde la suma es la habitual de los enteros.

Es trivial ver que los axiomas se cumplen:
\begin{enumerate}
\item La composición es asociativa: dadas $n, m, k : * \longrightarrow *$
(el único tipo de flechas que se puede componer, el único tipo de
flechas que hay) sabemos que
$n \circ (m \circ k) = n + (m + k) = (n + m) + k = (n \circ m) \circ k$.

\item Existe la identidad para cada objeto: solo existe un objeto y a su
identidad la llamaremos 0. Es trivial ver que $f \circ 0 = f$ y que
$0 \circ g = g$ en este contexto.
\end{enumerate}
En general esta construcción que acabamos de aplicar a $(\mathbb{Z}, +)$
se puede aplicar a cualquier monoide. Toda categoría con un solo
objeto se puede interpretar como un monoide (y viceversa):
la asociatividad de la
composición garantiza la asociatividad de la operación monoidal y
la existencia de la flecha identidad garantiza la existencia del
elemento neutro del monoide. En este sentido podemos considerar que
las categorías son una generalización de los monoides.

\paragraph{Categoría producto}
Dado un par de categorías $\C$ y $\D$ podemos construir la categoría
producto $\C\times\D$ de la siguiente forma:
\begin{itemize}
\item Los objetos serán de la forma $(C, D)$ donde $C$ es un objeto
  de $\C$ y $D$ es un objeto de $\D$.
\item Las flechas serán de la forma $(f, g) : \arr{(C, D)}{(C', D')}$
  donde $f: \arr{C}{C'}$ es una flecha de $\C$ y $g : \arr{D}{D'}$
  es una flecha de $\D$.
\item La composición actúa componente a componente
  $(f, g) \circ (f', g') = (f \circ f', g \circ g')$ siempre
  que $f$ y $f'$, y $g$ y $g'$ se puedan componer.
\end{itemize}

Las identidad del objeto $(C, D)$ es claramente la flecha
$(1_C, 1_D)$. Es trivial comprobar que $\C\times\D$ cumple
los axiomas de una categoría.

\section{Funtores}
\subsection{Definición}
De la misma manera que para los grupos se definen los homomorfismos
de grupos, para los anillos los homomorfismos de anillos y para los
espacios topológicos las funciones continuas, también
podemos asociar a las categorías una noción de morfismos que preserva
su estructura. A estos morfismos de categorías les llamaremos
\emph{funtores}. Un funtor $F$ de una categoría
$\C$ en una categoría $\D$ (que notaremos por $F : \arr{\C}{\D}$)
tendrá que llevar objetos de $\C$ en objetos de $\D$ y flechas
de $\C$ en flechas de $\D$ preservando la estructura de la
categoría en el siguiente sentido:
\begin{enumerate}
\item $F$ respeta los dominios y los codominios:
si $f : \arr{A}{B}$ una flecha
de la categoría $\mathcal{C}$ entonces
$F f : \arr{F A}{F B}$ es la correspondente flecha asociada
en $\D$.
\item $F$ preserva las identidades. Dicho de otra forma si $C$ es un
objeto de $\mathcal{C}$ entonces $F 1_C = 1_{F C}$.
\item $F$ respeta la composición: si tenemos $f: \arr{A}{B}$ y
$g : \arr{B}{C}$ tenemos que $F (g\circ f) = F g \circ F f$.
\end{enumerate}

Aunque la acción sobre los objetos
no determina por completo a un funtor, habrá ocasiones en las que
el lector podrá completar por sí mismo fácilmente la acción de este
sobre las flechas. En tales casos nos limitaremos a referirnos
al funtor describiendo su acción sobre los objetos.

La ley de la composición se puede interpretar como que $F$ lleva
diagramas conmutativos a diagramas conmutativos. \textcolor{red}{para hacer este comentario aquí tendrías que decir lo que  es un diagrama conmutativo ?¿?¿ Sí, diré qué es un diagrama conmutativo (antes de esta parte para que se entienda el comentario).}


\subsection{Ejemplos}
\paragraph{Funtores Identidad}
Para cada categoría $\C$ podemos definir el funtor identidad
$1_{\C} : \arr{C}{C}$ que deja invariante tanto a los objetos
como a las flechas. Comprobar las propiedades de los funtores
es trivial.

\paragraph{Funtores subyacentes}
Podemos considerar el funtor $U : \arr{\Grp}{\Set}$
que asigna a cada grupo su conjunto subyacente y a cada flecha la
aplicación entre los conjuntos subyacentes (cada homomorfismo de
grupos es también una aplicación entre ambos conjuntos). Es rutinario
comprobar que se respetan los axiomas de funtores. Existen también
funtores subyacentes desde la categoría de anillos, espacios topológicos
o retículos por ejemplo.

\paragraph{Grupos libres}
Podemos definir un funtor $F : \arr{\Set}{\Grp}$ de la siguiente
forma: a cada conjunto $X$ le asignamos el grupo libre sobre $X$
(al que llamaremos $F X$) y a
cada aplicación $f : \arr{X}{Y}$ entre conjuntos le asignamos el
único homomorfismo de grupos $F f: \arr{F X}{F Y}$ que extiende a $f$.
Comprobar que
$F$ es en efecto un funtor es sencillo.

\paragraph{Funtores $\Hom$}
Ya hemos dicho que para cada par de objetos $A, B$ de
una categoría $\C$ tenemos que
$\Hom(A, B)$ es un conjunto. Fijado un objeto $A$ de
$\C$ tenemos que $\Hom(A, -)$ nos permite asociar a cada
objeto $B$ de la categoría $\C$ un conjunto $\Hom(A, B)$.

Veamos que $\Hom(A, -)$ es un funtor
$\Hom(A, -) : \C \rightarrow \Set$. Ya conocemos la acción
sobre los objetos ahora tenemos que encontrar como actúa
el funtor sobre las flechas. Supongamos que tenemos
$f : \arr{B}{C}$ una flecha de $\C$. Definimos
$\Hom(A, f) : \arr{\Hom(A, B)}{\Hom(A, C)}$
($\Hom(A, f)$ es una aplicación entre los conjuntos
$\Hom(A, B)$ y $\Hom(A, C)$ es decir a cada función $\arr{A}{B}$
le asigna una función $\arr{A}{C}$) por
$$\Hom(A, f)(g) = f \circ g $$.

Probamos a modo de ejemplo
que se cumplen los axiomas de los funtores. En primer lugar
supongamos que $g : \arr{C}{D}$ y $h : \arr{D}{E}$ entonces tenemos
que probar
$$\Hom(A, h \circ g) = \Hom(A, h) \circ \Hom(A, g) :
\arr{\Hom(A, B)}{\Hom(A, C)}$$

Para probar tal cosa suponemos $f \in \Hom(A, B)$ y entonces:
\begin{multline*}
\Hom(A, h\circ g)(f) = (h \circ g) \circ f = h \circ (g \circ f) \\
= h \circ \Hom(A, g)(f) = \Hom(A, h)(\Hom(A, g)(f)) \\
= (\Hom(A, h) \circ \Hom(A, g))(f)
\end{multline*}

Y por tanto se comporta bien respecto a la composición. Veamos que se
comporta bien respecto a la identidad. Sea $f \in \Hom(A, B)$ entonces
$$\Hom(A, 1_B)(f) = 1_B \circ f = f$$ y por tanto
$\Hom(A, 1_B) = 1_{\Hom(A, B)}$.

Teniendo en cuenta $\Hom(A, -)$ se comporta bien respecto a la
composición y lleva identidades en identidades concluimos que es
un funtor.

\paragraph{Bifuntores}
Llamamos bifuntor a un funtor de
la forma $F : \arr{\C_1\times\C_2}{\D}$. Un ejemplo de
bifuntor sería $-\times - : \arr{\Set\times\Set}{\Set}$ que
a cada par de conjuntos le asigna su producto cartesiano.

%% \paragraph{Endofuntores en Haskell}
%% Es habitual que en lenguajes de programación fuertemente tipados exista
%% una noción de \emph{constructor de tipo} o \emph{generics}. Uno de los usos más
%% habituales de estos constructores de tipo son los contenedores. En
%% \verb~C++~ por ejemplo tenemos el constructor de tipo \verb~vector~, que no es
%% un tipo en si mismo pero que asigna a cada tipo un nuevo tipo. Por
%% ejemplo al tipo \verb~int~ le asigna el tipo \verb~vector<int>~ que representa
%% un vector con elementos de tipo \verb~int~.

%% En Haskell existen también constructores de tipos. Viéndolo desde el
%% punto de vista de \verb~Hask~ un constructor de tipos es una forma de asignar
%% objetos de la categoría a otros objetos de la categoría. Esto nos hace
%% pensar que un funtor de \verb~Hask~ a \verb~Hask~ (cuando un funtor va de una
%% categoría en sí misma le llamamos endofuntor) será un constructor de
%% tipos que además cumple una serie de propiedades adicionales.

%% Pondremos un ejemplo de endofuntor en \verb~Hask~. De forma similar a
%% \verb~vector~ podemos considerar el constructor de tipos \verb~[]~ de Haskell
%% que a cada tipo (por ejemplo \verb~Int~) le hace corresponder otro tipo
%% (en este caso \verb~[Int]~, una lista de enteros). Nos hace falta ahora
%% definir cómo a una función \verb~f :: a -> b~ se le asigna una función
%% de tipo \verb~[a] -> [b]~ de manera que se cumplan los axiomas
%% de los functores, pero eso es sencillo: extenderemos una función \verb~f~
%% sobre listas aplicando \verb~f~ elemento a elemento. Escribimos en haskell
%% la definición de esta extensión y escribimos algunos ejemplos:

%% \begin{verbatim}
%% map :: (a -> b) -> ([a] -> [b])
%% map f [] = []
%% map f (x:xs) = (f x):(map f xs)


%% add_1 :: Int -> Int
%% add_1 x = x + 1

%% map add_1 [] ; Esto da como resultado []
%% map add_1 [1,2,3,4,5] ; [2,3,4,5,6]

%% ; en general map f [x1, x2,x3...] = [f x1, f x2, f x3, ...]
%% \end{verbatim}

%% Veamos que se cumplen los axiomas de los funtores. En primer lugar
%% veamos que la identidad va a la identidad:

%% \begin{verbatim}
%% map id [] ; esto es []
%% map id [x1, ..., xn]; esto es [id x1, ..., id xn] = [x1, ..., xn]
%% \end{verbatim}

%% por lo que la función \verb~map id~ es la identidad. Veamos que la función
%% además respeta las composiciones:

%% \begin{verbatim}
%% ((map f) . (map g)) [x1, ..., xn] = map f (map g [x1, ..., xn])
%%                                   = map f [g x1, ..., g xn]
%%                                   = [ f (g x1), ..., f (g xn) ]
%%                                   = [ (f . g) x1, ..., (f . g) xn ]
%%                                   = map (f . g) [x1, ..., xn]
%% \end{verbatim}

\section{En programación}
\subsection{Categorías}
\paragraph{$\Hask$}
En el contexto del lenguaje de programación Haskell (aunque esta
construcción es análoga en otros lenguajes de programación
fuertemente tipados) se suele hablar de la categoría
$\Hask$  en la que los objetos son los tipos del lenguaje
(por ejemplo \verb~Int~, \verb~String~ o \verb~Double~) y
las flechas son las funciones entre esos tipos. Por ejemplo
la función \verb~length :: String -> Int~ vista en
$\Hask$ sería una flecha
$\texttt{length} \in \Hom_{\Hask}(\texttt{String}, \texttt{Int})$.
Como operador de composición tenemos la composición habitual de
funciones, que en haskell se nota con \texttt{.} (un punto) y se
define de la siguiente forma:
\begin{verbatim}
(.) :: (b -> c) -> (a -> b) -> (a -> c)
(.) f g argumentWithTypeA = f (g argumentWithTypeA)
\end{verbatim}
Además tenemos la función \texttt{id} que nos define las
flechas identidad en $\Hask$:

\begin{verbatim}
id :: a -> a
id x = x
\end{verbatim}
Nótese que aun teniendo esta colección de objetos, de flechas, la
operación de composición (que es asociativa) y las identidades tenemos
que $\Hask$ no es una categoría. Esto se debe
entre otras cosas a algunas
peculiaridades del comportamiento del
valor especial \texttt{undefined} de Haskell. Salvando el uso de
este valor, la teoría de categorías es un modelo ampliamente aceptado
para el estudio de $\Hask$. No presumimos en este trabajo
de que $\Hask$ cumpla bajo toda circunstancia los axiomas de las
categorías pero eso no nos impedirá utilizar la teoría de categorías
para analizar y razonar sobre construcciones hechas sobre Haskell
(siendo conscientes de la imperfección del modelo).
Una justificación de que $\Hask$ es indistinguible de una
categoría restringiéndose a un subconjunto del lenguaje lo encontramos
en \cite{fastandloose}.

A lo largo del trabajo veremos como especializando construcciones
categóricas a $\Hask$ obtendremos aplicaciones de la teoría
de la categorías a la programación.

\paragraph{\texttt{Pipes}}
\textcolor{red}{Hablar de este otro ejemplo de categorías}


\subsection{Funtores}
\paragraph{Endofuntores en $\Hask$}
Llamamos endofuntor a un funtor que va de una categoría $\C$ en sí misma.
Sabiendo esto podemos comenzar a entender en qué consiste un endofuntor
en $\Hask$: una forma de asignar tipos a otros tipos (los objetos
de $\Hask$) y transformar funciones en otras funciones (las flechas
de $\Hask$) de manera que se preserven algunas relaciones.

Es habitual encontrar mecanismos que permiten asignar tipos a otros tipos
en los lenguajes de programación usados hoy día. En \verb~C++~
tenemos como ejemplo concreto \verb~vector~ que a cada tipo
(por ejemplo \verb~int~, un entero)
le asigna otro tipo (\verb~vector<int>~, un vector de enteros). En
general las templates de \texttt{C++} permiten realizar este tipo de
construcciones. En \verb~java~ los Generics cumplen una función similar.

También es habitual en los lenguajes de programación modernos
que las funciones
sean \textit{ciudadanos de primera clase} (\textit{first class citizens}),
es decir, se puede operar sobre las funciones como se opera sobre
cualquier otro valor del lenguaje. En este contexto es natural que
surjan funciones que reciben funciones como parámetro y devuelven
otras funciones. Python es un ejemplo de lenguaje en el que
se encuentran estas \textit{higher order functions} (se usan
tan habitualmente que hasta se incorporó en el
lenguaje sintaxis específica para ellas).

En la biblioteca estándar de Haskell existe una \verb~typeclass~
(el mecanismo de polimorfismo de Haskell, que para los propósitos
de este trabajo podemos suponer similar a las interfaces de java)
que sirve para dotar de comportamiento funtorial a los constructores
de tipo que implementemos. La typeclass
se llama, convenientemente, \verb~Functor~ y se
define de la siguiente manera:

\begin{verbatim}
class Functor F where
  fmap :: (a -> b) -> (F a -> F b)
\end{verbatim}

Si queremos que nuestro constructor de tipo
sea un \verb~Functor~ tendremos
que implementar sobre él una función llamada \verb~fmap~ que reciba
una función de tipo \verb~a -> b~ y nos devuelva una función
de tipo \verb~F a -> F b~ donde \verb~F~ es nuestro constructor de tipo.

Veremos ejemplos a continuación que aclararán la situación pero si
tuviéramos que trazar un paralelismo con \texttt{C++} podríamos decir
que \verb~vector~ (que sería \verb~F~ del código en Haskell) sería una
instancia de la typeclass \verb~Functor~ si implementáramos una función
llamada \verb~fmap~ que recibe como parámetro una función que va
de un tipo cualquiera \verb~A~ a un tipo cualquiera \verb~B~ y
devuelve una función que va del tipo \verb~vector<A>~ (análogo a
\verb~F a~ en el código en Haskell) a \verb~vector<B>~.

Haskell no comprueba que tu implementación de \texttt{fmap}
verifica los axiomas de los funtores. Esa tarea se delega al
desarrollador. Los axiomas de los funtores en
haskell teniendo en cuenta los
operadores de composición \verb~.~ y la función identidad son:

\begin{verbatim}
; f :: a -> b
; g :: b -> c

fmap (g . f) = (fmap g) . (fmap f) :: F a -> F c

fmap id = id :: F a -> F a
\end{verbatim}

Proponemos algunos ejemplos de instancias de la typeclass
\verb~Functor~
que se encuentran en la biblioteca estándar
de haskell.

\paragraph{Maybe}
La definición de \verb~Maybe~ es la siguiente:

\begin{verbatim}
data Maybe a = Just a | Nothing
\end{verbatim}

Este tipo se usa constantemente en Haskell. Representa el resultado
de computaciones que podrían fallar o podrían no tener solución en
casos concretos. Podemos poner un ejemplo de
utilización de este tipo: \verb~head_safe~, una
función que devuelve el primer elemento de una lista:

\begin{verbatim}
head_safe :: [a] -> Maybe a
head_safe [] = Nothing
head_safe (x:xs) = Just x
\end{verbatim}

Decidimos que el tipo de retorno de \verb~head_safe~ sea \verb~Maybe a~ puesto
que este cómputo puede no tener solución en caso de que la lista no
tenga elementos. En otros lenguajes la función \verb~head~ lanzaría una
excepción si se le pasara una lista vacía, pero en Haskell se puede
codificar la naturaleza propensa a fallos
del resultado en el sistema de tipos.

Resulta que se puede dotar al constructor de tipos \verb~Maybe~ de
un comportamiento funtorial. Mostramos a continuación su
implementación de la typeclass \verb~Functor~

\begin{verbatim}
instance Functor Maybe where
  fmap f (Just x) = Just (f x)
  fmap f Nothing = Nothing
\end{verbatim}

Lo que hace esta función \verb~fmap~ es extender funciones de tipo
\verb~a -> b~ a una función de tipo \verb~Maybe a -> Maybe b~. Esta función
no hace nada si recibe un \verb~Nothing~ (representante del fallo
en el cómputo) y aplica la función al contenido del valor \verb~Just x~.

Podemos comprobar que se cumplen los axiomas de los funtores.

\begin{verbatim}
; veamos fmap id = id :: Maybe a -> Maybe a
; fmap id x = id x = x
; si x = (Just y) entonces
(fmap id) x = fmap id (Just y) = Just (id y) = (Just y) = x

; si x = Nothing
(fmap id) Nothing = Nothing

; veamos que se comporta bien con la composición.
; una vez más supongamos que x = (Just y)
(fmap (f . g)) x = fmap (f . g) (Just y)
                 = Just ( (f . g) y )
                 = Just (f (g y))
                 = (fmap f) (Just (g y))
                 = ( (fmap f) . (fmap g) ) (Just y)
                 = ( (fmap f) . (fmap g) ) x

; si x = Nothing
(fmap (f . g)) x = (fmap (f . g)) Nothing
                 = Nothing
                 = (fmap f) Nothing
                 = (fmap f) ((fmap g) Nothing)
                 = ((fmap f) . (fmap g)) Nothing
\end{verbatim}

\paragraph{Either}
La definición de \verb~Either~ es la siguiente:

\begin{verbatim}
data Either a b = Left a | Right b
\end{verbatim}

El tipo \verb~Either~ en haskell se usa para representar cómputos que pueden
devolver valores de dos tipos distintos. Un ejemplo muy habitual
de \verb~Either~ es representar cómputos que, al igual que \verb~Maybe~,
podrían ser erróneos, pero dando detalles sobre el error en caso
de error.

Proponemos un ejemplo artificial que aun así muestra para qué se
podría usar este tipo. Supongamos que tenemos un sistema con usuarios
registrados y en nuestra empresa queremos premiar la fidelidad
de nuestros usuarios en edad laboral. Supongamos además que tenemos
dos tipos de premio: uno para adultos jóvenes y otro para el resto
de personas en edad laboral.

Utilizando \verb~Maybe~ para resolver el problema nos quedaría un código
de la siguiente forma:

\begin{verbatim}
data PremiosTrabajadores = PremiosJovenes | PremiosMayores

dar_premio :: Int -> Maybe PremiosAdultos
dar_premio age
  | age < 16 = Nothing
  | 16 <= age < 40 = Just PremiosJovenes
  | 40 <= age < 65 = Just PremiosMayores
  | 65 <= age = Nothing
\end{verbatim}

Este código cumple su propósito de decirnos qué premio
le corresponde al usuario en caso de que efectivamente le
toque un premio. Sin embargo, lo que no devuelve la función
es el motivo por el que el cliente no es elegible para este.
Para conseguir que la función devuelva ese tipo de información
podemos usar \verb~Either~.

\begin{verbatim}
data PremiosTrabajadores = PremiosJovenes | PremiosMayores

dar_premio :: Int -> Either String PremiosAdultos
dar_premio edad
  | edad < 16 = Left "Demasiado joven para estar en edad laboral"
  | 16 <= edad < 40 = Right PremiosJovenes
  | 40 <= edad < 65 = Right PremiosMayores
  | 65 <= edad = Left "Demasiado mayor para estar en edad laboral"
\end{verbatim}

¿Es \verb~Either~ un funtor? La respuesta es que no, porque de entrada
\verb~Either~ es un constructor de tipo con dos parametros y para que
un constructor de tipo sea un funtor necesitamos que solo tenga
un parámetro. Entonces \verb~Either~ no es un funtor, pero resulta que
\verb~Either a~ donde \verb~a~ es algún (cualquier) tipo fijo de Haskell sí
es un funtor. Es decir si consideramos fijo el primer tipo
\verb~(Either a)~ es un constructor de tipos que admite un tipo
como parámetro y además se puede implementar una instancia de
\verb~Functor~ sobre él de la siguiente forma:

\begin{verbatim}
instance Functor (Either a) where
  fmap f (Left x) = Left x
  fmap f (Right x) = Right (f x)
\end{verbatim}

Esta instancia de \verb~Functor~ es similar a la de \verb~Maybe~: si el valor
es de los de \emph{error} no se hace nada con él. Si es de los valores
\emph{buenos} se transforma mediante la función \verb~f~. Veamos que efectivamente
esta instancia de \verb~Functor~ cumple con las leyes:

\begin{verbatim}
; la identidad va a la identidad:
; supongamos x = (Left y)

fmap id x = fmap id (Left y) = Left y = x

; supongamos x = Right y
fmap id x = fmap id (Right y) = Right (id y) = Right y = x

; probemos ahora que se lleva bien con la composición
fmap (f . g) (Left y) = Left y = (fmap f) (Left y)
                      = (fmap f) (fmap g (Left y))
                      = (fmap f) . (fmap g) (Left y)

fmap (f . g) (Right y) = Right ( (f . g) y )
                       = fmap f (Right (g y))
                       = (fmap f . fmap g) (Right y)
\end{verbatim}

Veamos un ejemplo de utilización de la instancia de \verb~Functor~ de
\verb~Either a~ siguiendo con el ejemplo que utilizamos antes. Imaginemos
que tenemos una función que asocia los distintos premios a sus títulos.
Por ejemplo:

\begin{verbatim}
titulos_premios :: PremiosTrabajadores -> String
titulos_premios PremioJovenes = "Semana de senderismo"
titulos_premios PremioMayores = "Cata de Vinos"
\end{verbatim}

Entonces si quisiéramos una función que a partir de la edad de
un usuario nos devolviera qué mensaje mostrarle en la interfaz
con respecto al premio podríamos hacer lo siguiente:

\begin{verbatim}
mensaje_premio :: Int -> String
mensaje_premio edad =
  case resultado of
    (Left mensajeError) -> "Error: " ++ mensajeError
    (Right tituloPremio) -> "Enhorabuena has conseguido una " ++ tituloPremio

  where
    resultado :: Either String String
    resultado = fmap titulos_premios (dar_premio edad)
\end{verbatim}

\paragraph{List}
\paragraph{Reader}
Definimos \texttt{Reader} de la siguiente forma:
\begin{verbatim}
data Reader a b = Reader (a -> b)
\end{verbatim}

Esta definición quiere decir
que un valor de \textit{tipo} \texttt{Reader a b} es
de la forma \texttt{Reader g} donde \texttt{g} es
una función que recibe un valor de tipo \texttt{a} como
parámetro y devuelve un valor de tipo \texttt{b}. De la misma manera
que hicimos con \texttt{Either} podemos fijar la primera variable
de tipo e implementar una instancia de \texttt{Functor} para
\texttt{(Reader a)}:

\begin{verbatim}
instance Functor (Reader a) where
  ; fmap :: (b -> c) -> (Reader a b) -> (Reader a c)
  fmap f (Reader g) = Reader (f . g)
\end{verbatim}

Podemos probar que esta instancia de \texttt{Functor} cumple
las leyes de los funtores:

\begin{verbatim}
; f :: b -> c
; g :: c -> d
; a  = (Reader h) :: (Reader a b) y entonces
; h :: a -> b

fmap (g . f) a = fmap (g. f) (Reader h)
               = Reader ( (g . f) . h)
               = Reader ( g . (f . h))
               = fmap g (Reader (f . h))
               = fmap g (fmap f (Reader h))
               = (fmap g . fmap f) (Reader h)
               = (fmap g . fmap f) a

; y por tanto fmap (g . f) = fmap g . fmap f
; en las mismas condiciones:

fmap id a = fmap id (Reader h)
          = Reader (id . h)
          = Reader h
          = a

;; y por tanto fmap id = id
\end{verbatim}

Veremos más adelante en el trabajo las aplicaciones
de \texttt{Reader}. Creemos también importante resaltar las similitudes
entre \texttt{Reader} y los funtores $\Hom$ descritos en los ejemplos
matemáticos de funtores. Formalizaremos esta relación
más adelante en el trabajo
cuando hablemos de objetos $\Hom$ internos.

\chapter{Construcciones elementales}
Dedicamos este capítulo al tratamiento de construcciones elementales
sobre categorías. En este capítulo podremos ver cómo un marco tan general
como el de la teoría de categorías permite realizar definiciones
\textit{Continuar introducción de construcciones elementales luego}.

\section{Elementos}
Cuando hablamos de conjuntos es común hablar de sus elementos. Muchas
definiciones que se hacen sobre conjuntos y aplicaciones entre
ellos se hacen en base a los elementos de los conjuntos involucrados.
Dos ejemplos de este tipo de definiciones serían las definiciones de
aplicación inyectiva y de producto cartesiano.

\begin{definition*}
Una aplicación $f : \arr{A}{B}$ es inyectiva si dados
$a, a' \in A$ tenemos que $f(a) = f(a') \implies a = a'$.
\end{definition*}

\begin{definition*}
Dados dos conjuntos $A$ y $B$ definimos su producto cartesiano
como:
$$A\times B = \{ (a, b) : a \in A \quad b \in B \}$$
\end{definition*}

Si intentamos trasladar las construcciones que hacemos sobre los
conjuntos y sus aplicaciones a categorías arbitrarias con sus objetos
y sus flechas tenemos que saber cómo trasladar la noción de elemento.

Si consideramos el conjunto con un solo elemento, al que llamaremos
$*$, y las aplicaciones que salen de él nos damos cuentas de que podemos
identificar las aplicaciones entre el conjunto $\uno = \{ 1 \}$
$A$ (estas aplicaciones son de la forma $1 \mapsto a \in A$)
arbitrario con los elementos (en el sentido de teoría de conjuntos)
de $A$. Dicho de otra forma $\Hom(*, A)$ y $A$ son ``\emph{lo mismo}"  como conjuntos.
pero $\Hom(*, A)$ está formado por flechas y eso es algo de lo que
sí podemos hablar dentro de una categoría. Sin embargo, para realizar
este procedimiento de identificar los elementos de un conjunto con
un conjunto de flechas de la categoría hemos tenido que acudir a un
objeto especial de la categoría de conjuntos: $\uno$. Este procedimiento
es algo que quizá no podamos realizar en otras categorías pero nos
motiva a hacer una definición más general de
\textit{elemento categórico}:

\begin{definition}
En una categoría $\C$ llamaremos elemento de un objeto $A$
a cualquier flecha $x : \arr{T}{A}$ (sea cual sea el objeto
$T$). De forma paralela a como se hace con conjuntos, utilizaremos
la notación $x \in^T A$. Diremos también que $x$ es un elemento
de $A$ definido sobre $T$.
\end{definition}

Con esta noción de elemento notamos que $f : \arr{A}{B}$ lleva
elementos de $A$ a elementos de $B$ mediante la composición: si
$x \in^T A$ entonces $f \circ x \in^T B$. Esto motiva que en lo que
sigue omitamos a veces el signo de composición ($f x$ en lugar
de $f \circ x$) si nuestra intención es
que algunas flechas (en este caso $x$) se interpreten
como elementos categóricos.

Veremos como esta noción nos permite llevar a teoría de categorías
algunos conceptos originados sobre conjuntos.

\section{Monomorfismos}
\begin{definition}
Consideremos una categoría $\C$ y una flecha $f : \arr{A}{B}$ de esta.
Diremos que $f$ es un monomorfismo si para cualquier objeto
$T$ de $\C$ y dados dos elementos de $A$
$x, y \in^T A$ tenemos que $fx = fy \implies x = y$.
\end{definition}

Nótese cómo esta definición es casi idéntica a la definición de
inyectividad sobre aplicaciones entre conjuntos. De hecho es sencillo
demostrar que la noción de monomorfismo sobre $\Set$ se corresponde
efectivamente con las aplicaciones inyectivas.


\paragraph{Ejemplos}
Esta noción se puede aplicar sobre otras categorías. En general
cuando consideramos categorías en la que los objetos son estructuras
matemáticas y las flechas son sus morfismos, los monomorfismos son
aquellas flechas tales que las aplicaciones entre conjuntos subyacentes
son inyectivas. Ejemplos de esto son las categorías
$\Grp$, $\Ring$, ...

\section{Isomorfismos}
\begin{definition}
Consideremos una categoría $\C$ y una flecha $f : \arr{A}{B}$ de
esta. Diremos que $f$ es un isomorfismo si $f$ es una biyección
entre los elementos de $A$ y los elementos de $B$ definidos sobre
$T$ para cualquier objeto $T$ de $\C$.
\end{definition}

Esta definición es similar a la definición habitual
de biyección sobre conjuntos pero generalizando los elementos de
los conjuntos a elementos de los objetos definidos sobre todos
los objetos
de la categoría.

Nótese que de este concepto podemos dar una definición equivalente
que no requiere del uso de elementos categóricos.

\begin{proposition}
Dada la categoría $\C$ y y una flecha $f: \arr{A}{B}$ tenemos que
$f$ es un isomorfismo sí y solo sí existe una flecha $g: \arr{B}{A}$
tal que $f \circ g = 1_B$ y $g \circ f = 1_A$.
\end{proposition}


Esta caracterización de isomorfismo nos permite ver rápidamente
que los isomorfismos de la categoría $\Set$ se corresponden con las
biyecciones y que los isomorfismos sobre $\Grp$, $\Ring$, $\Top$, ...
se corresponden con los isomorfismos de grupos,
isomorfismos de anillos y homeomorfismos de espacios topológicos.


\section{Productos}
También es posible trasladar la definición de producto de conjuntos
a una categoría arbitraria mediante la siguiente definición:

\begin{definition}
Sea $\C$ una categoría y $A$ y $B$ dos objetos de esta. Diremos que
existe el producto de $A$ y $B$ si
existe una terna $(A\times B, \pi_1, \pi_2)$
donde $A\times B$ es un objeto de $\C$ y
$\pi_1 : \arr{A\times B}{A}, \pi_2 : \arr{A\times B}{B}$ son dos flechas tales
que para cualquiera elementos $x : \arr{T}{A}$ de $A$ e
$y: \arr{T}{B}$ de $B$, existe un único elemento
$(x,y): \arr{T}{A\times B}$ de $A\times B$ tal que $\pi_1 (x,y)=x$ y $\pi_2(x,y)=y$.

Expresaremos esto diciendo que  el siguiente diagrama es conmutativo:
\begin{center}
\begin{tikzcd}
&\arrow{ld}[swap]{x} T \arrow[d, dotted, "{\exists! (x, y)}"] \arrow[rd, "y"] & \\
A & \arrow[l, "\pi_1"] A\times B  \arrow{r}[swap]{\pi_2} & B
\end{tikzcd}
\end{center}

\end{definition}

Debemos resaltar dos aspectos importantes de esta definición:
\begin{itemize}
\item El producto de dos objetos en una categoría dada no tiene por qué
      existir.
\item De existir un producto, este solo queda determinado salvo isomorfismo. De hecho si $(A\times B, \pi_1, \pi_2)$ es un producto de $A$ y $B$
y $\phi : \arr{P}{A\times B}$ es un isomorfismo entonces
$(P, \pi_1 \circ \phi , \pi_2 \circ \phi)$ es otro producto de $A$
y $B$.
\end{itemize}


\subsection{Ejemplos}
\paragraph{\Set}
Dados dos conjuntos $A$ y $B$ siempre existe el producto de estos
y además coincide (salvo el isomorfismo
que comentamos antes) con el producto cartesiano de ambos
junto a sus proyecciones canónicas. Desmotramos tal
afirmación a continuación:

Sean $x: \arr{T}{A}, y: \arr{T}{B}$ dos aplicaciones. Podemos
definir la aplicación $(x, y) :\arr{T}{A\times B}$ por
$(x, y)(t) = (x(t), y(t))$. Esta aplicación cumple en efecto que
$\pi_1 \circ (x, y) = x$ y $\pi_2 \circ y = y$. Pero además es la única
que cumple esto pues si tenemos otra aplicación
$f : \arr{T}{A\times B}$ cumpliendo $\pi_1 \circ f = x$ y
$\pi_2 \circ f = y$ solo nos basta recordar que podemos
escribir $f$ como:
$f(t)=(\pi_1\circ f(t), \pi_2\circ f(t))=(x(t), y(t))$.
y por tanto $f = (x, y)$.

\paragraph{Otras estructuras matemáticas}
A continuación mostramos ejemplos de productos
en categorías conocidas:

\begin{itemize}
\item En $\Grp$ el producto se corresponde con el
producto directo de grupos.
\item En $\Top$ se corresponde con el producto de espacios topológicos.
\item En $\Ring$ se corresponde con el producto de anillos.
\end{itemize}

\section{Objetos finales}
Al principio de la sección le dimos un papel especial al conjunto
de un solo elemento $\uno$. Nos gustaría caracterizar a ese objeto de
la categoría $\Set$ con alguna propiedad estrictamente categórica.
Esto es posible:

\begin{definition}
Sea $\C$ una categoría y $\uno$ un objeto. Diremos que $\uno$ es un objeto
final si dado cualquier objeto $T$ de la categoría existe un
único elemento de $\uno$ definido sobre $T$.
\end{definition}

Es sencillo ver cómo el conjunto $\uno$ se corresponde con el
objeto $\final$ de la categoría $\Set$.

\subsection{Ejemplos}
\paragraph{En $\Grp$}
El grupo trivial es el objeto final de la categoría $\Grp$.
Esto es sencillo de comprobar: dado cualquier grupo $G$ tenemos
que homomorfismo trivial $\arr{G}{\final}$ es el único morfismo
entre $G$ y $\final$.

\paragraph{En $\Ring$}
En la categoría de los anillos con unidad el objeto final es
también el anillo trivial.

\section{Dualidad}
El concepto de dualidad que se encuentra frecuentemente en las
matemáticas puede ser analizado de forma muy general desde el punto
de vista categórico. Para introducir esta noción presentamos a
continuación la definición de categoría opuesta.

\subsection{La categoría opuesta}
Dada una categoría $\C$ podemos construir una categoría
$\C^{op}$ en la que definimos $\Obj{\C^{op}} = \Obj{\C}$,
$\Hom_{\C^{op}}(A, B) = \Hom_{\C}(B, A)$ y una operación de
composición dada por $f \circ_{op} g = g \circ f$. Veamos
la construcción realizada describe en efecto una categoría:
\begin{itemize}
\item La composición es asociativa: sean
      $f \in \Hom_{\C^{op}}(A, B)$, $g \in \Hom_{\C^{op}}(B, C)$
      y $h \in \Hom_{\C^{op}}(C, D)$. Tenemos que
\begin{multline*}
(h \circ_{op} g) \circ_{op} f = (g \circ h) \circ_{op} f \\
                              = f \circ (g \circ h)
                              = (f \circ g) \circ h
                              = (g \circ_{op} f)  \circ h
                              = h \circ_{op} (g \circ_{op} f)
\end{multline*}
\item Existen las identidades. Las identidades siguen siendo
      las mismas flechas que son identidad en $\C$. La demostración
      es inmediata.
\end{itemize}

La categoría opuesta nos otorga una herramienta para reaprovechar
definiciones realizadas anteriormente. Veremos ejemplos de propiedades
duales en las siguientes secciones.

\textcolor{red}{aqui tienes que hablar de propiedades duales. Luego yo diría que llamaremos epimorfismos la propiedad dual de monomorfimo, coproducto a la dual de producto, final a la dual de inicial etc.  y luego diría que vas especificar estas propiedades más detalladamente }



\subsection{Epimorfismos}
En una sección anterior especificamos cuando una flecha
$f : \arr{A}{B}$ de una categoría $\C$ era un monomorfismo. La categoría
opuesta nos permite dualizar esta definición de la siguiente manera:

\begin{definition}
Dada una categoría $\C$ y una flecha $f: \arr{A}{B}$, diremos que
$f$ es un epimorfismo si y solo si $f$ es un monomorfismo en
$\C^{op}$.
\end{definition}

En este sentido decimos que las propiedades ``$f$ es un monomorfismo''
y ``$f$ es un epimorfismo'' son propiedades duales. Esta propiedad
se puede enunciar de forma sencilla también sin recurrir a
$\C^{op}$.

\begin{proposition}
$f$ es un epimorfismo si y solo si dado cualquier objeto $X$ y
cualquier par de flechas $g_1, g_2 : \arr{B}{X}$
tenemos que $g_1 \circ f = g_1 \circ f \implies g_1 = g_2$.
Decimos también en este caso que $f$ se puede cancelar por la
derecha.
\end{proposition}

\subsubsection{Ejemplos}
\paragraph{En \Set}
En la categoría de los conjuntos los epimorfismos coinciden con las
aplicaciones sobreyectivas.

\paragraph{Isomorfismos}
Todo isomorfismo es un epimorfismo y un monomorfismo. El recíproco
es cierto en algunas categorías (como en $\Set$ o $\Grp$) pero no
lo es en general.

\paragraph{En \Ring}
Consideremos la categoría de anillos con unidad (en la
que los objetos son anillos con unidad y
las flechas son homomorfismos de anillos). En esta categoría
que $f : \arr{R}{S}$ es equivalente también a que $f$ sea una aplicación
inyectiva.

Sin embargo consideremos la inclusión
$i : \arr{\Z}{\Q}$, un anillo $R$
y un par de aplicaciones de anillos $g_1, g_2 : \arr{\Q}{R}$. Supongamos
que $g_1 \circ i = g_2 \circ i$. Ahora  $\forall x \in \Q$
tenemos que $\exists a, b \in \Z : x = \frac{a}{b}$ y entonces:
\begin{multline*}
g_1(x) = g_1(\frac{a}{b}) = g_1(a)g_1(b)^{-1} \\
       = (g_1 \circ i)(a)(g_1\circ i)(b)^{-1}
       = (g_2\circ i)(a)(g_2\circ i)(b)^{-1}
       = g_2(a)g_2(b)^{-1} = g_2(x)
\end{multline*}

Con lo que $g_1=g_2$. Esto prueba que $i : \arr{\Z}{\Q}$ es epimorfismo
y sin embargo no es sobreyectiva como aplicación. Un isomorfismo en
la categoría de anillos se corresponde con un isomorfismo de anillos
(que en particular es una biyección de conjuntos) y por tanto
en la categoría de anillos se pueden encontrar flechas que son
monomorfismos y epimorfismos pero no son isomorfismos.

\subsection{Objetos finales}
Existe también una propiedad dual a la de ser objeto inicial. \textcolor{red}{toda propiedad tiene su dual, no queda bien decir que existe la propiedad dual de ser inicial, tendrías que decir que la propiedad dual de ser un objeto inicial es ser un objeto final o algo así}

\begin{definition}
Dado un objeto $T$ de una categoría $\C$ decimos que
$T$ es un objeto final si $T$ es un objeto inicial en la
categoría $\C^{op}$.
\end{definition}

Podemos probar una caracterización que no hace referencia a
$\C^{op}$: \textcolor{red}{yo no pondría esto como proposición, diría algo así: Si leemos lo que significa ser objeto final en la propia categoría $\C$ obtendríamos...}

\begin{proposition}
El objeto $T$ es final si y solo si para cualquier otro objeto
$X$ existe una sola flecha tiene a $X$ como dominio y a $T$ como
codominio.
\end{proposition}

\subsubsection{Ejemplos}
\paragraph{En \Set}
En $\Set$ el conjunto vacío $\emptyset$ es un objeto final.

\paragraph{En \Grp}

El objeto final en $\Grp$ es el grupo de un solo elemento. Notemos
que coincide con el objeto inicial. Cuando en una categoría $\C$ coinciden
el objeto inicial y el final $T$ llamamos a $T$ objeto nulo. El objeto
nulo nos permite definir una aplicación nula (\textit{explicar esto mejor}). \textcolor{red}{no tengo muy claro que nulo sea el nombre apropiado para esto. Yo siempre he usado cero}

Otro ejemplo de categoría en la que ocurre esto es en la de espacios vectoriales
sobre un cuerpo $K$.

\paragraph{En \Ring}
\textit{Probar que no existen objetos finales en la categoría de anillos} \textcolor{red}{que me dices del anillo con sólo un elemeno el $0$ que es al mismo tiempo cero y uno?¿?¿ lo que pasa en anillos es que no hay cero, el inicial y el terminal no coinciden}

\subsection{Coproducto}
Podemos repetir la definición de producto categórico sobre la categoría
opuesta y \textit{dando la vuelta} a las flechas para volver a la
categoría inicial $\C$ llegaríamos a la siguiente definición de lo que
llamaremos \textbf{coproducto}: \textcolor{red}{de nuevo yo diría aquí algo así: Coproducto es la propiedad dual de producto, si leemos esta propiedad en la propia categoría $\C$ en lugar de en $\C^{op}$ tendríamaos}

\begin{definition}
Sea $\C$ una categoría y $A$ y $B$ dos objetos de esta. Diremos que
existe el coproducto de $A$ y $B$
si existe una terna $(A+B, i_1, i_2)$
donde $P$ es un objeto de $\C$ y
$i_1 : \arr{P}{A}, i_2 : \arr{P}{B}$ son dos flechas tales
que para cualquier objeto $Y$ y sendas flechas $f_1 : \arr{A}{Y}$,
$f_2 : \arr{B}{Y}$ existe un morfismo
$f : \arr{P}{Y}$ tal que el siguiente diagrama es conmutativo:
\begin{center}
\begin{tikzcd}
&  Y  \arrow[rd, "f_2"]& \\
A \arrow[ur, "f_1"] \arrow{r}[swap]{i_1} &
 \arrow[u, dotted, "\exists! f"] P
 & B \arrow[l, "i_2"]
\end{tikzcd}
\end{center}
\end{definition}

\textit{Retocar esta sección para que sea más consistente con las otras}
\textit{Hablar de que es único salvo isomorfismo}. \textcolor{red}{Quizás deberías decir que la propiedad dual de isomorfismo es de nuevo isomorfismo, esto hace que el carácter \emph{único salvo isomorfismo} que tienen las construcciones definidas por  una \emph{propiedad universal} también lo tengan las construcciones definidas por sus propiedades duales.....}


\subsection{Funtores $\Hom(-, A)$}
\textcolor{red}{Creo que le das demasiadas vueltas a esto. Diría algo así; Dado un objeto $A$ de una categoría $\C$ definimos el funtor $Hom_\C(-,A)$ como $$Hom_\C(-,A)=Hom_{\C^{op}}(A,-):\arr{\C^{op}}{\Set}$$ Si queremos observar el efecto de este funtor sobre $\C$ tenemos que para cada flecha $\arr{C'}{C}$....}



Sea $\C$ una categoría y $A$ un objeto de esta.
Sea $f \in \Hom_{\C^{op}}(C, C')$. Definimos:

$$\Hom_{\C}(f, A) :
  \arr{
    \Hom_{\C}(C, A)
  }{
    \Hom_{\C}(C', A)
  }$$
$$\Hom_{\C}(f, A)(g) = g \circ f$$

(recordando que $g : \arr{C}{A}$ y $f : \arr{C'}{C}$ como
flechas de $\C$ y por tanto se pueden componer)

Se cumple además que:\textcolor{red}{esto es evidente}

\begin{itemize}
\item $\Hom_{\C}(1_C, A)(g) = g \circ 1_C = g$ para toda flecha
  $g\in \Hom_{\C}(C, A)$ y por tanto
  $\Hom_{\C}(1_C, A) = 1_{\Hom_{\C}(C, A)}$ para todo objeto $C$
  de $\C^{op}$.

\item Sean $f \in \Hom_{\C^{op}}(C, D)$ y $g : \Hom_{\C^{op}}(D, E)$
  entonces
  \begin{multline*}
    $$\Hom_{\C}(g \circ_{op} f, A)(h) = h \circ (g \circ_{op} f) \\
      h \circ (f \circ g) = (h \circ f) \circ g = \Hom_{\C}(g, A)(h \circ f) \\
      = \Hom_{\C}(g, A)(\Hom_{\C}(f, A)(h))$$
  \end{multline*}
  para toda $h \in \Hom_{\C}(C, A)$ y por tanto
  $\Hom_{\C}(g \circ_{op} f, A) = \Hom_{\C}(g, A) \circ \Hom_{\C}(f, A)$.
\end{itemize}

Esto nos muestra que $\Hom_\C(-, A)$ es un funtor
$\Hom_\C(-, A) : \arr{\C^{op}}{\Set}$, lo que no resulta
nada sorprendente teniendo en cuenta que
$\Hom_\C(-, A) = \Hom_{\C^{op}}(A, -)$ (que ya vimos que
era un funtor) tanto sobre los objetos
como sobre las flechas. Más aun, se puede demostrar
que $\Hom_\C(-, -) : \arr{\C^{op} \times \C}{\Set}$ es un bifuntor.

Cuando queremos hacer énfasis a que un funtor está actuando sobre
la categoría opuesta a una particular, decimos que es un funtor
contravariante. Por ejemplo, diríamos que $\Hom(-, A)$ es un funtor
contravariante \textcolor{red}{desde la categoría $\C$}.

\section{En programación}
A lo largo de esta sección aplicamos las construcciones que hemos visto
a la categoría \texttt{Hask}.

\subsection{Objetos iniciales y finales}
\texttt{Hask} tiene tanto objetos iniciales como objetos finales. El
objeto final es el tipo con un solo valor, al que se le suele llamar
\texttt{()} (pronunciado \texttt{Unit}). \texttt{()} es un tipo
cuyo único valor es \texttt{()} y para cada tipo \texttt{a} podemos
definir la función:
\begin{verbatim}
constUnit :: a -> ()
constUnit x = ()
\end{verbatim}
Por otro lado tenemos que el objeto inicial de \texttt{Hask} es el tipo
llamado \texttt{Void}, que no tiene ningún valor. En la biblioteca
estándar se encuentra una función llamada \texttt{absurd}
que tiene
como tipo \texttt{Void -> a} y es la única función con este tipo.
En cualquier caso, la función no puede ser llamada puesto que no
hay ningún valor con el que llamarla. La situación tanto como para
el objeto inicial como para el final es muy similar a la que se daba
en \Set.

\subsection{Productos}
\texttt{Hask} tiene productos. Dado dos tipos \texttt{A} y
\texttt{B} podemos construir el tipo \texttt{(A, B)} que tiene
como valores pares de valores de \texttt{A} y de \texttt{B}.
Las proyecciones se implementan en la biblioteca estándar de
Haskell de la siguiente forma:

\begin{verbatim}
fst :: (a, b) -> a
fst (x, y) = x

snd :: (a, b) -> b
fst (x, y) = y
\end{verbatim}

Si tenemos un par de funciones \cod{f1 :: X -> A} y
\cod{f2 :: X -> B} podemos definir la función \texttt{f}:
\begin{verbatim}
f :: X -> (A, B)
f x = (f1 x, f2 x)
\end{verbatim}

\subsection{Coproductos}
En la sección sobre funtores introdujimos el tipo \cod{Either}:
\begin{verbatim}
data Either a b = Left a | Right b
\end{verbatim}
Resulta que \cod{Either} es el coproducto en \cod{Hask}. Si
\cod{A} y \cod{B} son dos tipos de haskell entonces su coproducto
es \cod{Either A B} y las inyecciones canónicas son por un lado
\cod{Left :: A -> Either A B} y por otro lado
\cod{Right :: B -> Either A B}. Si tenemos un par de funciones
\cod{f1 :: A -> Y} y \cod{f2 :: B -> Y} tenemos que la única
función \cod{f :: Either A B -> Y} que se lleva bien con las inyecciones
es:
\begin{verbatim}
f :: Either A B -> Y
f (Left a) = f1 a
f (Right b) = f2 b
\end{verbatim}

Veamos que se lleva bien con las inyecciones:
\begin{verbatim}
(f . Left) a = f (Left a) = f1 a
;; por tanto f . Left = f1

(f . Right) b = f (Right b) = f2 b
;; por tanto f . Right = f2
\end{verbatim}


\chapter{Transformaciones Naturales y el lema de Yoneda}
\section{Transformaciones Naturales}
En el primer capítulo introdujimos los funtores como los morfismos
entre las categorías. Ahora ha llegado el momento de dar un paso más
e introducir las transformaciones naturales como morfismos entre
funtores. Esta es una de las nociones que motivó la creación de la
teoría de categorías: se encuentran ejemplos de esta en múltiples
ramas de las matemáticas.

Procedemos con la definición:

\begin{definition}
  Dados dos funtores $F: \arr{\C}{\D}, G : \arr{\C}{\D}$ decimos
  que $\lambda : \nat{F}{G}$ es una transformación natural si $\lambda$
  asigna a cada objeto $C$ de $\C$ una flecha
  $\lambda_C : \arr{FC}{GC}$ de manera que para
  cualquier flecha $g : \arr{C}{C'}$ el siguiente diagrama
  es conmutativo:
\begin{center}
    \begin{tikzcd}
      F C \arrow{r}{\lambda_C} \arrow{d}[swap]{F g} & GC \arrow{d}{G g} \\
      F C' \arrow{r}{\lambda_{C'}} & G C'
    \end{tikzcd}
  \end{center}
\end{definition}

Los siguientes resultados nos permitirán considerar categorías en las
que los objetos son funtores entre dos categorías $\C$ y $\D$ y las
flechas son las transformaciones naturales entre los funtores:

\begin{proposition}
  \begin{enumerate}
  \item Dado cualquier funtor $F : \arr{\C}{\D}$ podemos definir
    la transformación natural $1_F : \nat{F}{F}$ dada por
    ${(1_F)}_C = 1_{(FC)}$
    para cada objeto $C$ de $\C$.
  \item Podemos componer transformaciones naturales de la siguiente
    forma: dados funtores $F, G, H : \arr{\C}{\D}$ y transformaciones
    naturales $\lambda : \nat{F}{G}$, $\sigma : \nat{G}{H}$
    definimos la transformación natural
    $\sigma \circ \lambda : \nat{F}{H}$ por
    $(\sigma\circ\lambda)_C = \sigma_C \circ \lambda_C$.
  \item La composición de transformaciones naturales
    es asociativa en el siguiente sentido: dados
    $F \xrightarrow{\lambda} G \xrightarrow{\sigma} H \xrightarrow{\tau}I$
    tenemos que $(\tau \circ \sigma) \circ \lambda = \tau \circ (\sigma \lambda)$
  \item Dado cualquier par de funtores $F,G : \arr{\C}{\D}$ y
    transformaciones naturales $\tau : \arr{F}{G}$,
    $\sigma : \arr{G}{F}$ tenemos que $1_F \circ \sigma = \sigma$
    y $\tau \circ 1_F = \tau$.
  \end{enumerate}

  En definitiva los funtores $\arr{\C}{\D}$ y las transformaciones
  naturales entre estos funtores forman una categoría. A esta categoría
  la llamaremos $\funccat{\C}{\D}$.
\end{proposition}


\begin{definition}
  Dados dos funtores $F, G : \arr{\C}{\D}$ diremos que la transformación
  natural $\lambda : \nat{F}{G}$ es un isomorfismo natural si
  $\lambda_C : \arr{FC}{GC}$ es un isomorfismo para todo objeto $C$
  de $\C$.
\end{definition}

Nótese que los isomorfismos naturales son precisamente los isomorfismos
en las categorías de funtores.

\subsection{Ejemplos}
\paragraph{El doble dual}
Consideremos la categoría de espacios vectoriales sobre un cuerpo
$K$ a la que llamaremos $\VectK$. Podemos considerar los endofuntores
identidad $1_{\VectK} : \arr{\VectK}{\VectK}$ y  doble dual:
$(-)^{**} : \arr{\VectK}{\VectK}$. El último de estos funtores actúa
así sobre los morfismos:

$$(-)^{**} : \arr{\Hom(V, W)}{\Hom(V^{**}, W^{**})}$$
$$f^{**}(g)(\phi) = g(\phi \circ f)$$

Para cada espacio vectorial sobre $K$ $V$ \textcolor{red}{uff} podemos además
definir un monomorfismo de espacios vectoriales
 (isomorfismo
en el caso finito-dimensional) en su doble dual:

\begin{equation*}
i_V : \arr{V}{V^{**}}
\end{equation*}
\begin{equation*}
i_V(v)(\phi) = \phi(v)
\end{equation*}

Pues $i$ es una transformación natural entre $1_K$ y
$(-)^{**}$. Más aun, si nos restringimos al caso finito dimensional,
debido a que todas las componentes de la transformación natural
son isomorfismos, decimos que ambos funtores son naturalmente isomorfos. \textcolor{red}{este paraqraph no está muy bien redactado, además naturalmente isomorfos e isomorfos en la categoría de funtores ?¿?¿? a estas alturas no se?¿?? y también la restricción a finito dimensionales uff?¿?¿ estás cambiando los dominios de definición....?¿?¿}

\paragraph{Abelianización de grupos}
Dado un grupo $G$ definimos su abelianización $G^{ab}$ como
$G^{ab} = \frac{G}{[G, G]}$. Llamemos $\pi_G : \arr{G}{G^{ab}}$
a la proyección sobre el cociente. Para cualquier morfismo
de grupos $f : \arr{G}{H}$ tenemos que la aplicación
$\pi_H \circ f$ contiene en su nucleo a $[G, G]$ (ya que
$H^{ab}$ es abeliano y un morfismo de un grupo a un grupo abeliano
lleva los conmutadores al 0) y por tanto factoriza a través
de $G^{ab}$ permitiéndonos definir una aplicación
$f^{ab} : \arr{G^{ab}}{H^{ab}}$ y comprobar
que $(-)^{ab}$ es un funtor. La proyección de $\pi$ es una transformación
natural $\pi : \arr{ 1_{\Grp} }{ (-)^{ab} }$. \textcolor{red}{también tienes que redactar esto un poco mejor}

\section{Lema de Yoneda}
El lema de Yoneda es uno de los primeros resultados
inesperados que surgen a raíz del
estudio de la teoría de categorías. Una
de las consecuencias más importantes de este resultado es
el encaje de Yoneda, que nos permite
ver una categoría $\C$ arbitraria dentro de la categoría de
funtores $\funccat{\C^{op}}{\Set}$. Necesitaremos algunos resultados
previos para poder enunciar el lema.

A lo largo de las siguientes
secciones usaremos la notación $\Nat(F, G)$ en lugar de
$\Hom_{\funccat{\C}{\D}}(F, G)$ para referirnos al conjunto de
transformaciones naturales entre dos funtores
$F, G: \arr{\C}{\D}$. Probamos la funtorialidad de uno de los
ingredientes principales del lema de Yoneda.
\begin{proposition*}
  Dadas una categoría $\C$ y un objeto $A$ de esta definimos:
  $$L : \arr{\C\times\funccat{\C}{\Set}}{\Set}$$
  $$L(A, F) = \Nat(\Hom(A, -), F)$$
  $L$ es un funtor.
\end{proposition*}
\begin{proof}
  Llamando
  $H: \arr{\C^{op}}{\funccat{\C}{\Set}}$ al funtor dado por
  $H(A) = \Hom_\C(A, -)$
  podemos notar que nuestro funtor $L$ es precisamente
  $$L :\C\times\funccat{\C}{\Set}
  \xrightarrow{H^{op}\times 1_{\funccat{\C}{\Set}}}
  (\funccat{C}{\Set})^{op} \times \funccat{C}{\Set}
  \xrightarrow{\Nat} \Set$$
\end{proof}
Probamos la funtorialidad
del otro ingrediente principal del lema de Yoneda.
\begin{proposition*}
  Sea $\C$ una categoría y $Ev(-, -) : \arr{\C\times\funccat{\C}{\Set}}{\Set}$
  dado por $Ev(C, F) = FC$. $Ev$ es un funtor y lo llamamos
  \textbf{funtor de evaluación}.
\end{proposition*}
\begin{proof}
  Sean $\sigma : \arr{F}{G}$ y $\tau : \arr{G}{H}$ transformaciones naturales entre funtores
  $F, G, H : \arr{\C}{\Set}$. Sean además $f : \arr{C}{D}$,
  $g : \arr{D}{E}$ flechas de $\C$.
  Definimos la acción de $Ev$ sobre las flechas de
  $\arr{\C\times\funccat{\C}{\Set}}{\Set}$ como
  $Ev(f, \sigma) = \sigma_D \circ Ff : \arr{FC}{GD}$. Veamos que $Ev$ se comporta bien
  respecto a las composiciones:
  \begin{align*}
    Ev(g \circ f, \tau\circ\sigma) &
     = (\tau \circ \sigma)_E\circ F(g\circ f)
     = \tau_E\circ\sigma_E\circ Fg \circ Ff \\
     & \stackrel{*}{=}
  \tau_E\circ Gg \circ\sigma_D \circ Ff
  = Ev(g, \tau) \circ Ev(f, \sigma)
  \end{align*}

  Donde la igualdad $(*)$ se deduce de la
  naturalidad de $\sigma : \nat{F}{G}$.
  Que $Ev(1_C, 1_F) = 1_{FC}$ se prueba de forma sencilla,
  confirmando que $Ev : \arr{\C\times\funccat{\C}{\Set}}{\Set}$ es un funtor
\end{proof}
Ya estamos preparados para enunciar el lema.
\begin{theorem}[Lema de Yoneda]
  Sea $\C$ una categoría. Entonces los funtores $L$ y $Ev$ son naturalmente isomorfos.
  En particular, para cada objeto $A$ de $\C$ y cada funtor $F : \arr{\C}{\Set}$
  se tiene una biyección natural
  $$\Nat(\Hom(A, -), F) \cong FA$$
\end{theorem}
\begin{proof}
  Definimos la aplicación
  $$\phi_{A, F} : \arr{\Nat(\Hom(A, -), F)}{FA}$$
  $$\phi_{A, F}(\tau) = \tau_A (1_A)$$
  (recordamos que $\tau_A$ por ser $\tau$ transformación natural entre $\Hom(A, -)$ y $F$ es
  induce una aplicación de conjuntos $\tau_A : \arr{\Hom(A, A)}{F A}$).

  Veamos que $\phi_{A, F}$ es una biyección. En primer lugar veremos que es sobreyectiva. Sea $x \in FA$
  Definimos la transformación natural $\lambda_x : \nat{\Hom(A, -)}{F}$ por
  $$(\lambda_x)_C : \arr{\Hom(A, C)}{F C}$$
  $$(\lambda_x)_C(f) = Ff(x)$$

  Comprobar que $\lambda_x$ es en efecto una transformación natural es sencillo. Pero
  $\phi_{A, F}(\lambda_x) = (\lambda_x)_A (1_A) = F1_A(x) = x$, luego $\phi_{A, F}$
  es sobreyectiva.

  Veamos ahora que $\phi_{A, F}$ es inyectiva. Supongamos que tenemos
  dos transformaciones naturales $\tau, \tau' : \nat{\Hom(A, -)}{F}$
  tales que $\phi_{A, F}(\tau) = \phi_{A, F}(\tau')$. Para todo $f \in \Hom(A, C)$
  la naturalidad de $\tau$ nos garantiza la conmutatividad
  del siguiente diagrama:
  \begin{center}
    \begin{tikzcd}
      \Hom(A, A) \arrow{d}[swap]{\Hom(A, f)} \arrow{r}{\tau_A}
            & FA \arrow{d}{Ff} \\

      \Hom(A, C) \arrow{r}{\tau_C} & F C
    \end{tikzcd}
  \end{center}


  Es decir $\tau_C \circ \Hom(A, f) = F f \circ \tau_A$. Aplicando
  esta flecha sobre el valor $1_A$ tenemos que:
$$\tau_C(f) = \tau_C(f \circ 1_A) = \tau_C(\Hom(A, f)(1_A)) = F f(\tau_A(1_A)) = Ff(\phi_{A, F}(\tau))$$
  De la misma manera se puede obtener que $\tau'_C(f)  = Ff(\phi_{A, F}(\tau'))$, pero entonces para cualquier objeto $C$ y para cualquier flecha $f \in \Hom(A, C)$
  se cumple:
  $$\tau_C(f) = Ff (\phi_{A,F}(\tau)) = Ff(\phi_{A, F}(\tau')) = \tau'_C(f)$$
  con lo que $\tau = \tau'$ y $\phi_{A, F}$ es biyectiva.
  Merece la pena apuntar cuál es la aplicación inversa de $\phi_{A, F}$.
  Unas cuentas sencillas nos permiten comprobar que:
  $$\phi_{A,F}^{-1} : \arr{FA}{\Nat(\Hom(A, -), F)}$$
  $$\phi_{A,F}^{-1}(a)_C(f) = (Ff)(a)$$

  Probar que $\phi_{A, F}$ es natural en $A$ y $F$ es lo que queda para ver
  que $\phi : \nat{L}{Ev}$ es un isomorfismo natural. La demostración es
  rutinaria y no se incluye en el presente texto.
\end{proof}

El lema de Yoneda se puede interpretar en términos de elementos
categóricos. Si tenemos una categoría $\C$, un objeto $A$ de esta
y un funtor $F : \arr{\C}{\Set}$, el lema de Yoneda nos permite
decir que los elementos categóricos de $F$ definidos sobre
el funtor $\Hom(A, -)$ son esencialmente los mismos que los
elementos del conjunto $F A$.

Obtenemos un importante corolario del Lema de Yoneda cuando tomamos
$F = \Hom(B, -)$ con $B$ otro objeto de $\C$.

\begin{theorem}
  Dada una categoría $\C$
  el funtor $T : \arr{C^{op}}{\funccat{\C}{\Set}}$ dado por
  $T(A) = \Hom_\C(A, -)$ es fiel y pleno.
\end{theorem}
\begin{proof}
  $T$ actúa sobre flechas
  $f \in \Hom_{\C^{op}}(A, B)$
  de la siguiente manera:
  $$T(f) = \Hom_\C(f, -) : \nat{T(A) = \Hom(A, -)}{\Hom(B, -) = T(B)}$$
  $$T(f)_C(g) = g \circ f$$

  Repitiendo la demostración del lema de Yoneda con
  $F = \Hom(B, -)$ vemos que la aplicación inducida por
  $T$ sobre los conjuntos $\Hom_{\C^{op}}(A, B)$ es precisamente la
  biyección:
  $$\phi_{A, F}^{-1} : \arr{
    \Hom_{\C^{op}}(A, B) = FA
  }{
    \Nat(\Hom(A, -), F) = \Nat(T(A), T(B))
  }$$
  que definimos anteriormente.
  Luego $T$ es inyectiva y sobreyectiva sobre los conjuntos
  $\Hom$ y
  por tanto es un funtor fiel y pleno.
\end{proof}

Aplicando esta misma proposición sobre la categoría $\C^{op}$ llegamos
al resultado dual:

\begin{theorem}
  Dada una categoría $\C$, el funtor que asigna a cada objeto
  $A$ de la categoría el funtor $\Hom(-, A)$ es un funtor
  $\arr{\C}{\funccat{\C^{op}}{\Set}}$ fiel y pleno.
\end{theorem}

Los funtores descritos en los dos últimos teoremas son conocidos
como \textbf{encajes de Yoneda}.

Los encajes de Yoneda son una herramienta fundamental para la teoría
de categorías. Nos permiten ver la categoría $\C$ dentro de la
categoría de funtores $\funccat{\C^{op}}{\Set}$. En este sentido
la categoría $\funccat{\C^{op}}{\Set}$ se puede ver como una extensión
de la categoría $\C$, identificando cada objeto $A$ de $\C$
con el funtor $\Hom(-, A)$ y cada flecha $f : \arr{A}{B}$ con la
transformación natural asociada
$\Hom(-, f) : \nat{\Hom(-, A)}{\Hom(-, B)}$. El hecho
de que el encaje de
Yoneda sea fiel y pleno nos dice que esta identificación es buena
en el sentido de que las flechas entre objetos $A$ y $B$ de $\C$
son \textit{las mismas} que las flechas (transformaciones naturales)
entre los funtores
$\Hom(-, A)$ y $\Hom(-, B)$

Cuando en el capítulo 2 introdujimos la noción
de elemento categórico subrayamos que dados $x \in^T A$ y
$f : \arr{A}{B}$ entonces $f \circ x \in^T B$.
Que $f$ lleve elementos de $A$ en elementos de $B$
motivó que utilizáramos la notación $f(x)$ en lugar de $f \circ x$.
Esta notación podría llevarnos a interpretar $f$ como una familia
de aplicaciones
$f_T : \arr{\Hom(T, A)}{\Hom(T, B)}$ para objeto $T$ de $\C$,
pero esta familia es precisamente
la imagen de $f$ a través del encaje de Yoneda: $f_T = \Hom(-, f)_T$.
Recíprocamente,
dada cualquier familia de aplicaciones
$g_T : \arr{\Hom(T, A)}{\Hom(T, B)}$ natural en $T$, en el sentido
de que para cualquier $h : \arr{S}{T}$ se dé
$g_T(x) \circ h = g_S(x \circ h)$, tendríamos que $g$ es una transformación
natural $g: \nat{\Hom(-, A)}{\Hom(-, B)}$ y el
encaje de Yoneda nos permitiría realizar la identificación $g : \arr{A}{B}$.

Todo esto nos permite decir que dar una flecha $f: \arr{A}{B}$ en una categoría
arbitraria $\C$ es esencialmente lo mismo que dar una aplicación que lleva
elementos de $A$ a elementos de $B$ de una forma coherente (la condición
de naturalidad).

\section{Programación}
\subsection{Transformaciones Naturales}
Para construir una transformación natural entre dos endofuntores de
$\Hask$ \cod{F} y \cod{G} necesitamos, en primer lugar, una función de tipo
\cod{FA -> GA} para cada tipo \cod{A} del lenguaje.

Usamos el sistema de polimorfismo de Haskell para suplir nuestra primera
necesidad. El lenguaje nos permite definir una función de
tipo \cod{F a -> G a} (usamos \cod{a} minúscula
cuando nos referimos a una función polimórfica en la variable
\cod{a}) que, especializando a cada tipo del lenguaje, será
la componente en cada objeto de $\Hask$ de nuestra transformación natural.

Un ejemplo de este tipo de función polimórfica sería:
\begin{minted}{haskell}
discardLeft :: Either b a -> Maybe a
discardLeft (Left stringQueNoQuiero) = Nothing
discardLeft (Right v) = Just v
\end{minted}
(recordemos que es \cod{Either b}, y no \cod{Either} a secas,
la instancia de la clase \cod{Functor})

En segundo lugar, necesitamos que se cumplan las condiciones de
naturalidad. Podríamos emplear nuestro tiempo en demostrar que
la función \cod{discardLeft} que hemos definido cumple
tales condiciones pero eso no
será necesario: el sistema de polimorfismo de
Haskell (\textit{parametric polymorphism}) nos lo garantiza.
Haskell nos impone la restricción de que
una función polimórfica en la variable de tipo \cod{a}
esté implementada
con la misma expresión independientemente de futuras
especializaciones de \cod{a}. Por ejemplo, en Haskell no es posible
escribir una función polimórfica de esta forma:
\begin{minted}{haskell}
polimorfica :: a -> a
polimorfica v = v

polimorfica :: Int -> Int
polimorfica n = n + 1
\end{minted}
Esta restricción, junto con otras características
del polimorfismo paramétrico, permite comprobar que cualquier
función de tipo \cod{F a->G a} (donde, insistimos, la \cod{a}
minúscula significa que se trata de una función polimórfica
en el tipo \cod{a})
con \cod{F} y \cod{G} funtores es
una transformación natural \cite{theorems-for-free}.

Mostramos a continuación algunos ejemplos de transformaciones naturales
que podemos escribir entre los funtores de Haskell que hemos visto
anteriormente.

\subsubsection{Ejemplos}
\paragraph{Entre \cod{Maybe} y \cod{Either}}
La función \cod{discardLeft :: Either b a -> Maybe a} que hemos definido
antes es un ejemplo de transformación natural. La interpretación
que dimos de estos tipos en el primer capítulo fue la siguiente.
\begin{itemize}
\item Los valores de la forma \cod{Nothing :: Maybe a} representan
  situaciones de error durante un cómputo. Los valores de la forma
  \cod{Just v :: Maybe a} representan un cómputo exitoso que tiene
  al valor \cod{v :: a} como valor de retorno.
\item Los valores de la forma \cod{Left error :: Either b a} representan
  situaciones de error durante un cómputo y además aportan más información
  sobre el error (por ejemplo si la variable \cod{b} estuviera
  especializada a \cod{String} podríamos tener un valor de error
  que fuera \cod{Left "No fue posible conectar con el servidor"}. Los
  valores de la forma \cod{Right v :: Either b a} representan cómputos
  exitosos con valor de retorno \cod{v :: a}.
\end{itemize}

En este sentido podemos interpretar la función \cod{discardLeft} como
una aplicación que recibe un cómputo posiblemente fallido y descarta
la información asociada al fallo en caso de que la haya. En caso de
computación exitosa deja invariante su resultado.

Podemos escribir una transformación natural en el sentido inverso
fijando un valor de tipo \cod{b}. Por ejemplo:
\begin{minted}{haskell}
maybeToRight :: b -> Maybe a -> Either b a
maybeToRight defaultError Nothing = Left defaultError
maybeToRight defaultError (Just v) = Right v
\end{minted}

Y \cod{maybeToRight} no es una transformación natural pero para
cualquier valor \cod{x :: B} tenemos que
\cod{maybeToRight x :: Maybe a -> Either B a} sí que lo es.

\paragraph{Transformaciones naturales al funtor constante}
En Haskell podemos definir el funtor constante de la siguiente
manera:
\begin{minted}{haskell}
data Const b a = Const b

instance Functor (Const b) where
  fmap f (Const b) = Const b
\end{minted}

Este funtor es un funtor trivial porque lleva a todos los tipos
esencialmente al mismo tipo y además lleva todas las funciones,
también esencialmente, a la
identidad.
Supongamos que tenemos una función polimórfica de tipo
\cod{f :: F a -> B} donde \cod{B} es un tipo concreto. Esta función
se puede extender a otra función
\cod{fConst :: F a -> Const B a} de la siguiente forma:
\begin{minted}{haskell}
fConst :: F a -> Const B a
fConst x = Const (f x)
\end{minted}
El tipo de \cod{fConst} (es de la forma \cod{F a -> G a} donde tanto
\cod{F} como \cod{G} son funtores) nos garantiza que es una transformación
natural. Entonces cumple las condiciones de naturalidad. Especificando
la condición de naturalidad con la sintaxis de Haskell llegamos a
que para cualquier función \cod{g :: C -> C'} (donde \cod{C}
y \cod{C'} son dos tipos fijos pero arbitrarios):
\begin{minted}{haskell}
fConst . fmap g = fmap g . fConst
\end{minted}
donde el \cod{fmap} que aparece en el miembro
de la izquierda es el de la instancia
de \cod{Functor} de \cod{F} y el de
la derecha es el de \cod{Const B}.
Pero vimos anteriormente que la implementación de \cod{fmap}
de \cod{Const b} es básicamente la identidad.
Esto nos lleva a comprobar que:
\cod{f . fmap g = f}. Es decir, dada cualquier función de tipo
\cod{f :: F a -> B} los valores \cod{f} son invariantes a través
de aplicaciones de la forma \cod{fmap g}. Ponemos dos ejemplos
que aclararán esta proposición:

\begin{itemize}
\item Consideramos la función \cod{length :: [a] -> Int}. En este
  caso lo que estamos diciendo es que dada una lista
  \cod{l :: [A]} y cualquier función \cod{h :: A -> B} tenemos que
\begin{minted}{haskell}
length (fmap h l) = length l
\end{minted}
   Es decir, la longitud de
  una lista no cambia a través de funciones de la forma
  \cod{fmap h}.
\item Consideramos ahora la función
\begin{minted}{haskell}
isNothing :: Maybe a -> Bool
isNothing Nothing = True
isNothing (Just x) = False
\end{minted}
En este caso dado un valor \cod{m :: Maybe A} y cualquier función
\cod{h :: A -> B} tenemos que:
\begin{minted}{haskell}
isNothing (fmap f m) = isNothing m
\end{minted}
Es decir, que las funciones de la forma \cod{fmap h} no pueden llevar
un valor \cod{Nothing :: Maybe C} a uno
\cod{Just x :: Maybe D} ni viceversa.
\end{itemize}

\paragraph{Transformaciones naturales a través de funtores}
Realizamos ahora la construcción de las transformaciones naturales
a través de los funtores en Haskell. Utilizaremos el tipo
\cod{Composition f g a} que usamos en uno de los ejemplos de
funtores.
\begin{minted}{haskell}
natApplication :: (Functor f, Functor g, Functor h) =>
  (f a -> g a) -> Composition h f a -> Composition h g a
natApplication tau (Composition hfa) =
  Composition (fmap tau hfa)
\end{minted}
Esta construcción es la que nos permite dada una transformación
natural \cod{tau :: F a -> G a} y un funtor \cod{H} obtener una transformación
natural \cod{natApplication tau :: H F a -> H G a}.
La transformación natural \cod{F H a -> G H a} se obtiene simplemente
restringiendo \cod{tau} a tipos de la forma \cod{H a}.

Veamos un ejemplo de uso de \cod{natApplication}. Especializaremos
\cod{f} al funtor \cod{Either String}, \cod{g} al funtor
\cod{Maybe} y \cod{h} al funtor \cod{[]}. La transformación natural
que usaremos será \cod{discardLeft :: Either String a -> Maybe a}.
A continuación ponemos algunos ejemplos de cómo funciona
\cod{natApplication discardLeft}.
\begin{minted}{haskell}
-- natApplication discardLeft :: [Either String a] -> [Maybe a]
natApplication
    discardLeft
    (Composition [Left "error", Right 5, Left "test", Right 6])

-- resultado: [Nothing, Just 5, Nothing, Just 6]

natApplication
    discardLeft
    (Composition [Right "gg", Right "hh", Left "hh"])
-- resultado [Just "gg", Just "hh", Nothing ]
\end{minted}

%% \subsection{Lema de Yoneda}
%% Aplicando el mismo razonamiento que utilizamos con el funtor
%% $\Hom(A, -)$ al funtor de haskell \cod{Reader a}
%% (esto lo podemos hacer porque $\Hask$ es una categoría cartesianamente
%% cerrada y \cod{Reader a b} es el objeto $\Hom$ interno, aunque
%% no profundizaremos en esto) podemos decir que
%% el tipo \cod{F A} donde \cod{F} es un funtor es isomorfo
%% al tipo de las transformaciones naturales \cod{Reader A a -> F a}.
%% Este isomorfismo lo podemos mostrar en haskell de la siguiente forma:
%% \begin{minted}{haskell}
%% data Yoneda f a = Yoneda (forall b. Reader a b -> f b)

%% iso :: (Functor f) => f a -> Yoneda f a
%% iso fa = Yoneda natTrans
%%   where
%%     natTrans :: Reader a b -> f b
%%     natTrans (Reader aToB) = fmap aToB fa

%% iso_inv :: (Functor f) => Yoneda f a -> f a
%% iso_inv (Yoneda tau) = tau (Reader id)

%% -- iso . iso_inv = id
%% \end{minted}

%% En vista de este isomorfismo podemos implementar una instancia
%% de \cod{Functor} para \cod{Yoneda f}:

\chapter{Adjunciones y Mónadas}
\section{Adjunciones}
De camino a nuestro objetivo de definir el concepto de mónada
pasaremos \textcolor{red}{pasaremos?¿??¿ comenzamos?¿?¿} por la definición de adjunción. Como veremos en las próximas
secciones las mónadas y las adjunciones están estrechamente relacionadas.
Esto no debe hacernos pensar que las mónadas son el único motivo
por el que las adjunciones son importante en matemáticas. Las adjunciones
han sido identificadas por algunos autores como uno de los conceptos
clave de la teoría de categorías. McLane resume la importancia de
esta noción diciendo que ``hay funtores adjuntos por todas partes''
\cite{ariseeverywhere}.

Empezaremos por precisar qué es una adjunción.

\subsection{Definición}
\begin{definition}
  Una adjunción entre categorías $\C$ y $\D$ es un par de funtores
  $F : \arr{\C}{\D}$ y $G: \arr{\D}{\C}$ tales que los funtores
  $\Hom_\D(F-, -) : \arr{\C^{op}\times\D}{\Set}$ y
  $\Hom_\C(-, G-) : \arr{\C^{op}\times\D}{\Set}$ son naturalmente isomorfos.

  Diremos que $F$ es el
  adjunto por la izquierda de $G$ y que $G$ es el adjunto
  por la derecha de $F$. Notaremos esta relación entre ambos funtores
  como $F \adjoint G$.
\end{definition}

Nótese que para cualquier
objeto $C$ de $\C$ y cualquier objeto
$D$ de $\D$ un par de funtores adjuntos $F \adjoint G$ nos permite
establecer biyecciones entre los conjuntos:

$$\Hom_\D(FC, D) \cong \Hom_\C(C, GD)$$

Sea $\phi_{C, D} : \arr{\Hom_\D(FC, D)}{\Hom_\C(C, GD)}$ una tal
familia de biyecciones. Comprobar que $\phi$ es natural se reduce
a, por un lado, comprobar que el siguiente diagrama en $\D$:

\begin{center}
  \begin{tikzcd}
    FC \arrow{r}{f} & D \arrow{r}{g} & D'
  \end{tikzcd}
\end{center}

Se traslada al siguiente diagrama conmutativo en $\C$:

\begin{center}
  \begin{tikzcd}
    C \arrow[bend right=50, swap]{rr}{\phi_{C, D'}(g \circ f)}
      \arrow{r}{\phi_{C,D}(f)} & GD \arrow{r}{G g} & GD'
  \end{tikzcd}
\end{center}

Y por otro lado, que el siguiente diagrama en $\C$:

\begin{center}
  \begin{tikzcd}
    C' \arrow{r}{h} & C \arrow{r}{f} & GD
  \end{tikzcd}
\end{center}

se traslada al siguiente diagrama conmutativo en $\D$

\begin{center}
  \begin{tikzcd}
    FC' \arrow[bend right=50, swap]{rr}{\phi^{-1}_{C', D}(f \circ h)}
        \arrow{r}{Fh} & FC \arrow{r}{\phi^{-1}_{C, D}(f)} & D
  \end{tikzcd}
\end{center}

\subsection{Unidad y counidad}
Dadas categorías $\C$ y $\D$ y funtores adjuntos $F \adjoint G$ es
posible construir dos transformaciones naturales
$\eta : \nat{1_\C}{GF}$ y $\epsilon: \nat{FG}{1_\D}$, a las que llamaremos
respectivamente la \textbf{unidad} y la \textbf{counidad} de la
adjunción. \textcolor{red}{creo que hay muchos las...mejor sería:..que llamaremos..la... o a las que llamaremos unidad y...}

Sea $\phi : \nat{\Hom(F-, -)}{\Hom(-, G-)}$ el isomorfismo natural que
nos viene dado por la adjunción. Damos la transformación $\eta$
por sus componentes
$\eta_C = \phi_{C, FC}(1_{FC}) : \arr{C}{GFC}$.

Veamos que en efecto es natural. Sea $f : \arr{C}{C'}$
una flecha de $\C$. El diagrama cuya conmutatividad
tenemos que verificar es el siguiente:

\begin{center}
  \begin{tikzcd}
    C \arrow{d}{f}
      \arrow{r}{\eta_C} & GFC \arrow{d}{GFf} \\

    C' \arrow{r}{\eta_{C'}} & GFC'
  \end{tikzcd}
\end{center}

Es decir tenemos que comprobar que
$$\phi_{C', FC'}(1_{FC'}) \circ f = \eta_{C'} \circ f = GFf \circ \eta_C
  = GFf \circ \phi_{C, FC}(1_{FC})$$

Pero teniendo cuenta que $\phi$ es natural tenemos que el diagrama

\begin{center}
  \begin{tikzcd}
    FC \arrow{r}{1_{FC}} & FC \arrow{r}{Ff} & FC'
  \end{tikzcd}
\end{center}

Lo podemos transponer en el siguiente diagrama conmutativo:

\begin{center}
  \begin{tikzcd}
    C
    \arrow[bend right=50, swap]{rr}{\phi_{C, FC'}(f \circ 1_C) = \phi_{C, FC'}(Ff) }
      \arrow{r}{\eta_C}
            & GFC \arrow{r}{GFf} & GFC'
  \end{tikzcd}
\end{center}

Luego $GFf \circ \eta_C = \phi_{C, FC'}(Ff)$ pero a su vez:


\begin{center}
  \begin{tikzcd}
    C \arrow{r}{f} & C'
       \arrow{r}{\eta_{C'}} & GFC'
  \end{tikzcd}
\end{center}

Lo podemos transponer en:

\begin{center}
  \begin{tikzcd}
    FC \arrow[bend right=50, swap]{rr}{\phi_{C, FC'}^{-1}(\eta_{C'} \circ f)}
       \arrow{r}{F f} & FC' \arrow{r}{1_{FC'}} & FC'
  \end{tikzcd}
\end{center}

Este último diagrama muestra que
$Ff = \phi^{-1}_{C, FC'}(\eta_{C'} \circ f)$.


Y por tanto
$$GFf \circ \eta_C = \phi_{C, FC'}(Ff) = \phi_{C, FC'}
  (\phi^{-1}_{C,FC'}(\eta_{C'} \circ f)) = \eta_{C'} \circ f$$

Y la naturalidad de $\eta : \nat{1_\C}{GF}$ queda probada.

\subsection{Ejemplos}
\paragraph{Producto}
Sea $A$ un conjunto. Veamos que los funtores
$-\times A : \arr{\Set}{\Set}$ y
$\Hom(A, -) : \arr{\Set}{\Set}$ conforman una adjunción.
Necesitamos encontrar un isomorfismo natural
$$\phi : \nat{\Hom(-\times A, -)}{\Hom(-, \Hom(A, -))}$$.
\textcolor{red}{el punto lo tienes que poner antes del doble dolar para que no pase esto}Para ello definimos
$$\phi_{B, Z} : \arr{\Hom(B\times A, Z)}{\Hom(B, \Hom(A, Z))}$$
$$\phi_{B, Z}(f)(b)(a) = f(b, a)$$

Veamos que $\phi_{B, Z}$ es una biyección.
En primer lugar vemos que es sobreyectiva.
Sea $f' \in \Hom(B, \Hom(A, Z))$. Definimos la aplicación
$f : \arr{B\times A}{Z}$ dada por $f(b, a) = f'(b)(a), \,  \forall b\in B \, , a \in A$. Aplicando
la definición de $\phi_{B, Z}$
tenemos $\phi_{B,Z}(f)(b)(a) = f(b, a) = f'(b)(a)$ también para
todo $b \in B$ y para todo $a \in A$, y por tanto
$\phi_{B,Z}(f) = f'$. Veamos ahora que $\phi$ es inyectiva. Supongamos
que tenemos $f_1, f_2 : \arr{B\times A}{Z}$ tales que
$\phi_{B, Z}(f_1) = \phi_{B, Z}(f_2)$ pero entonces tenemos que
$f_1(b, a) = \phi_{B, Z}(f_1)(b)(a) = \phi_{B,Z}(f_2)(b)(a) = f_2(b, a)$
$\forall (b, a) \in B\times A$
y por tanto $f_1 = f_2$. \textcolor{red}{este párrafo es un poco farragoso, la puntuación está regular además mezclas para todo símbolo con para todo texto, $\forall b\in B a\in A $ y $\forall (b,a)\in B\times A$ uff....}

Comprobamos ahora las condiciones de naturalidad. En primer lugar dado
el diagrama:
\begin{center}
  \begin{tikzcd}
    B \arrow{r}{f} & B' \arrow{r}{g} & \Hom(A, Z)
  \end{tikzcd}
\end{center}
Hemos de ver que el siguiente diagrama es conmutativo:
\begin{center}
  \begin{tikzcd}
    {B\times A}
      \arrow[bend right=50, swap]{rr}{\phi^{-1}_{B, Z}(g \circ f)}
      \arrow{r}{f\times A} & B'\times A \arrow{r}{\phi^{-1}_{B', Z}(g)} & Z
  \end{tikzcd}
\end{center}
Y esto equivale a probar
$$\phi_{B,  Z}(g \circ f) = \phi_{B', Z}(g) \circ (f\times A)$$
donde $f \times A$ está definida por
$$f\times A: \arr{B\times A}{B'\times A}$$
$$(f\times A)(b, a) = (f(b), a)$$
Aplicando sobre un valor cualquiera
$(b, a) \in B\times A$ tenemos que:
$$\phi_{B, Z}^{-1}(g \circ f)(b, a) = (g\circ f)(b)(a) = g(f(b))(a)$$
y
$$(\phi_{B', Z}^{-1}(g) \circ (f\times A))(b, a) =
   \phi_{B', Z}^{-1}(g)(f(b), a) = g(f(b))(a)$$
En segundo lugar dado el diagrama:
\begin{center}
  \begin{tikzcd}
    B\times A \arrow{r}{f} & Z \arrow{r}{g} & Z'
  \end{tikzcd}
\end{center}
Hemos de ver que el siguiene diagrama es conmutativo:
\begin{center}
  \begin{tikzcd}
    B \arrow[bend right=50, swap]{rr}{\phi_{B, Z'}(g \circ f)}
      \arrow{r}{\phi_{B, Z}(f)} & \Hom(A, Z) \arrow{r}{\Hom(A, g)} & \Hom(A, Z')
  \end{tikzcd}
\end{center}
Es decir, $\phi_{B, Z'}(g \circ f) = \Hom(A, g) \circ \phi_{B, Z}(f)$. Para
cualesquiera
$b \in B, a \in A$ tenemos que:
$$\phi_{B, Z'}(g \circ f)(b)(a) = (g\circ f)(b, a) = g(f(b, a))$$
y
\begin{multline*}
(\Hom(A, g)\circ\phi_{B, Z}(f))(b)(a) = (\Hom(A, g)(\phi_{B, Z}(f)(b))(a)\\
  = (g \circ \phi_{B,Z}(f)(b))(a) = g(f(b, a))
\end{multline*}
Y por tanto $\phi$ es una transformación natural.

Esta adjunción nos dice que dar una aplicación que va de $B\times A$ a $Z$
es esencialmente lo mismo que dar una aplicación que vaya de $B$ al conjunto
de aplicaciones entre $A$ y $Z$.
\paragraph{Grupos libres}
Los funtores de conjunto subyacente $U : \arr{\Grp}{\Set}$ y el funtor
$F : \arr{\Set}{\Grp}$ que asigna a cada conjunto su grupo libre asociado
son adjuntos $F \adjoint U$. \textcolor{red}{Ufff. yo simplificaría esto: Los funtores $U : \arr{\Grp}{\Set}$, conjunto subyacente, y $F : \arr{\Set}{\Grp}$ , grupo libre, son adjuntos  $F \adjoint U$.}

La unidad $\eta : \nat{1_\Set}{UF}$ es la inclusión \textcolor{red}{canónica}
$\eta_X : X \subseteq UFX$ de un conjunto en el
grupo libre generado por él. La counidad
$\mu : \nat{FU}{1_\Grp}$ determina para cada grupo $G$
el homomorfismo de grupos $\mu_G : \arr{FUG}{G}$
tal que a cada palabra
$g_1g_2\ldots g_n \in FUG$ le asigna el valor $\prod_{i=1}^n g_i \in G$.


\section{Mónadas}
Comenzamos con el estudio de una de las nociones nacidas
en teoría de categorías que son
más importantes para nuestro trabajo: las mónadas.


\subsection{Definición}
\begin{definition}
  Una mónada sobre una categoría $\C$
  es una terna $(T, \eta, \mu)$ donde $T : \arr{\C}{\C}$ es un
  endofuntor y $\eta : \nat{1_\C}{T}, \, \mu: \nat{T^2}{T}$ son dos
  transformaciones naturales tales que los siguientes diagramas
  conmutan:
  \begin{center}
    \begin{equation}\label{diagrama:mu}
    \begin{tikzcd}
      T^3 \arrow{d}[swap]{\mu T} \arrow{r}{T\mu} & T^2 \arrow{d}{\mu} \\
      T^2 \arrow{r}{\mu} & T
    \end{tikzcd}
    \end{equation}

    \begin{equation}\label{diagrama:eta}
      \begin{tikzcd}
        T \ar[equal]{dr} \arrow{r}{T\eta} & T^2 \arrow{d}{\mu} & T \arrow{l}[swap]{\eta T} \ar[equal]{ld} \\
        & T &
      \end{tikzcd}
    \end{equation}
  \end{center}
\end{definition}
Estudiar algunos ejemplos será útil para comprender mejor
los distintos elementos que forman parte
de la definición.

\subsubsection{Ejemplos}
\paragraph{Monoides}
Sea $M$ un monoide. Definimos:
\begin{itemize}
\item El funtor $T : \arr{\Set}{\Set}$
  como $T(X) = M\times X$
\item $\eta_X : \arr{X}{T(X)}$ definida
por $\eta_X(x) = (e, x)$ (con $e \in M$ el elemento neutro)
\item $\mu_X : \arr{T^2(X)}{T(X)}$ definida por
$\mu_X(m, n, x) = (m\cdot n, x)$
\end{itemize}

Veamos que $(T, \eta, \mu)$ es una mónada.
\textit{A lo largo de este ejemplo nos tomaremos la libertad de identificar $T^2X$ con $M\times M\times X$ en lugar con $M \times (M\times X)$ con la intención de simplificar la notación}.
En primer lugar
tenemos que ver que para cualquier conjunto $X$ el siguiente
diagrama es conmutativo (diagrama \eqref{diagrama:mu} de la definición
de mónada):
\begin{center}
  \begin{tikzcd}
    M\times M \times M \times X \arrow{r}{T\mu_X}
           \arrow{d}{\mu_{TX}} & M\times M\times X \arrow{d}{\mu_X} \\
    M\times M\times X \arrow{r}{\mu_X} & M \times X
  \end{tikzcd}
\end{center}
Es decir $\mu_X \circ \mu_{TX} = \mu_X \circ T\mu_X$. Para comprobar
esta igualdad tomamos $(m, n, k, x) \in T^3(X)$.
Por un lado:
$$(\mu_X \circ \mu_{TX})(m, n, k, x) = \mu_X(\mu_{TX}(m, n, k, x)
  = \mu_X(m\cdot n, k, x) = ((m\cdot n)\cdot k, x)$$
Por otro lado:
$$(\mu_X \circ T\mu_X)(m, n, k, x) = \mu_X(T\mu_X(m, n, k, x))
  = \mu_X(m, n\cdot k, x) = (m \cdot (n\cdot k), x)$$
y por la asociatividad del producto en $M$ tenemos que
el diagrama es conmutativo.
Queda ver que para cualquier conjunto $X$ el diagrama
(\ref{diagrama:eta} de la definición de mónada)
\begin{center}
  \begin{tikzcd}
    T(X) \ar[equal]{dr}
         \arrow{r}{\eta_{TX}} & T^2(X) \arrow{d}{\mu_X} & \arrow{l}[swap]{T\eta_X} \ar[equal]{ld} T(X) \\
    & T(X) &
  \end{tikzcd}
\end{center}
es conmutativo, es decir, $\mu_X \circ T\eta_X = 1_{TX} = \mu_X \circ \eta_{TX}$. Podemos
comprobar esta igualdad basta ver que para cualquier $(m, x) \in TX$:
$$\mu_X(T\eta_X(m, x)) = \mu_X(m, e, x) = (m \cdot e, x) = (m, x)$$
$$\mu_X(\eta_{TX}(m, x)) = \mu_X(e, m, x) = (e\cdot m, x) = (m, x)$$

\paragraph{Sumas}
Sea $\C$ una categoría con coproductos finitos. Dado un objeto
$C$ de $\C$ de definimos:

\begin{itemize}
\item El endofuntor $T: \arr{\C}{\C}$ dado por $T(X) = C + X$
\item La transformación natural $\eta : \nat{1_\C}{T}$ dada por
  $\eta_X : \arr{X}{C+X}$ la inyección canónica en el coproducto.
\item La transformación natural $\mu: \nat{T^2}{T}$ dada por
  $\mu_X : \arr{C + (C + X)}{C + X}$ \textit{Definir}
\end{itemize}

La terna $(T, \eta, \mu)$ es una mónada.


\paragraph{Mónada $\Hom(A, -)$}
Sea $A$ un conjunto. Definimos:

\begin{itemize}
\item El endofuntor $T = \Hom(A, -) : \arr{\Set}{\Set}$
\item La transformación natural $\eta : \nat{1_\Set}{T}$ dada
  por $\eta_X : \arr{X}{\Hom(A, X)}$ con $\eta_X(x)(a) = x$.
\item La transformación natural $\mu: \nat{T^2}{T}$ dada por
  $\mu_X : \arr{\Hom(A, \Hom(A, X))}{\Hom(A, X)}$ con
  $\mu_X(f)(a) = f(a)(a)$
\end{itemize}

La terna $(T, \eta, \mu)$ es una mónada.

\subsection{Las adjunciones dan lugar a mónadas}
\begin{theorem}
Sean $\C$ y $\D$ dos categorías y $F : \arr{\C}{\D}$ y $G: \arr{\D}{\C}$ un
par de funtores adjuntos $F \adjoint G$ con unidad
$\eta : \nat{1_{\C}}{GF}$ y counidad $\epsilon: \nat{FG}{1_\D}$. La terna
$(GF, \eta, G\epsilon F)$ es una mónada.
\end{theorem}
\begin{proof}
  El diagrama \eqref{diagrama:mu} de la definición de mónada es en este
  caso el diagrama:
  \begin{center}
    \begin{tikzcd}
      GF \arrow{r}{GF\eta} & GFGF \arrow{d}{G\epsilon F} & GFGF \arrow{l}[swap]{\eta GF} \\
      & GF \ar[equal]{ur} \ar[equal]{ul} &
    \end{tikzcd}
  \end{center}
  Es decir, hay que probar que:
  $$G\epsilon F \circ GF\eta = 1_{GF} = G\epsilon F \circ \eta GF$$
  Probaremos la igualdad de la derecha.
  Sea $C$ un objeto de $\C$. La componente en
  $C$ de la transformación natural del
  miembro de la derecha de la ecuación
  es $G\epsilon_{FC} \circ \eta_{GFC}$.
  Por la condición \eqref{diagrama:adjuncion1} de la definición de
  adjunción tenemos que el siguiente diagrama en $\D$:
  \begin{center}
    \begin{tikzcd}
      FGFC \arrow{r}{1_{FGFC}}
      & FGFC \arrow{r}{\epsilon_{FC}}
      & FC
    \end{tikzcd}
  \end{center}
  Lo podemos trasladar al siguiente diagrama conmutativo en $\C$:
  \begin{center}
    \begin{tikzcd}
      GFC \arrow{r}{\eta_{GFC}}
          \arrow[bend right=50, swap]{rr}{\phi_{GFC, FC}(\epsilon_{FC})}
      & GFGFC \arrow{r}{G\epsilon_{FC}}
      & GFC
    \end{tikzcd}
  \end{center}
  Y entonces $$\phi_{GFC, FC}(\epsilon_{FC}) = G\epsilon_{FC} \circ \eta_{GFC}$$
  Pero por la definición de $\epsilon_{FC}$ tenemos que
  $$1_{GFC} = \phi_{GFC, FC}(\epsilon_{FC}) = G\epsilon_{FC} \circ \eta_{GFC}$$
  Con lo que
  $G\epsilon F \circ \eta GF = 1_{GF}$. La prueba de que $G\epsilon F \circ GF\eta = 1_{GF}$ es análoga.


  Queda probar la condición \eqref{diagrama:eta} de la definición
  de mónada. El diagrama de \eqref{diagrama:eta} es en este caso:
  \begin{center}
    \begin{tikzcd}
      GFGFGF \arrow{d}{GFG\epsilon F}
             \arrow{r}{G\epsilon F GF}
      & GFGF \arrow{d}{G \epsilon F} \\
      GFGF \arrow{r}{G\epsilon F} & GFGF
    \end{tikzcd}
  \end{center}
  Y la conmutatividad de este diagrama es
  equivalente a la del siguiente:
  \begin{center}
    \begin{tikzcd}
      FGFG \arrow{d}{FG \epsilon}
           \arrow{r}{\epsilon FG}
      & FG \arrow{d}{\epsilon} \\
      FG \arrow{r}{\epsilon}& 1_\D
    \end{tikzcd}
  \end{center}
  Pero este diagrama en cualquier objeto $D$ de la categoría $\D$ es
  \begin{center}
    \begin{tikzcd}
      FGFG D \arrow{d}{FG\epsilon_D}
      \arrow{r}{\epsilon_{FGD}}
      & FG \arrow{d}{\epsilon_D} \\
      FG \arrow{r}{\epsilon_D} & D
    \end{tikzcd}
  \end{center}
  que es precisamente la condición de naturalidad de
  $\epsilon : \nat{FG}{1_\D}$ aplicada a la flecha
  $\epsilon_D : \arr{FGD}{D}$.
\end{proof}


\subsection{Categoría de Kleisli de una mónada}
Hemos visto que todo par de funtores adjuntos inducen una mónada.
Es natural plantearse si el recíproco es cierto: ¿es toda mónada
la composición de un par de funtores adjuntos?
La respuesta es afirmativa.

\begin{theorem}
  Sea $\C$ una categoría y $(T, \eta, \mu)$ una mónada
  sobre $\C$. Entonces existe una categoría $\D$ y un par
  de funtores $F : \arr{\C}{\D}$, $G : \arr{\D}{\C}$ tales
  que $F \adjoint G$ y además:
  \begin{enumerate}
  \item $T = GF$
  \item $\eta$ es la unidad de la adjunción
  \item $\mu = F\epsilon G$ donde $\epsilon: \nat{FG}{1_\D}$ es
    la counidad de la adjunción.
  \end{enumerate}
\end{theorem}
\begin{proof}[\textbf{Demostración. } Categoría de Kleisli]
  Probamos el resultado paso por paso.
  \begin{enumerate}
  \item Construcción de $\D$. Definimos $\Obj{\D} = \Obj{\C}$ y
    $\Hom_\D(A, B) = \Hom_\C(A, TB)$.
    Dadas dos flechas $f \in \Hom_\D(A, B)$ y
    $g \in \Hom_\D(B, C)$ definimos la composición $g \circ_\D f \in \Hom_\D(A, C)$
    como la flecha
    $$g \circ_\D f :
    A \xrightarrow{f} TB \xrightarrow{Tg} T^2 C \xrightarrow{\mu_C} T C
    $$
    Veamos que esta operación de composición
    cumple los axiomas:
    \begin{itemize}
    \item Es asociativa. Sean $f : \Hom_\D(A, B)$, $g : \Hom_\D(B, C)$
      y $h : \Hom_\D(C, D)$ entonces
      $$(h \circ_{\D} g) \circ_\D f = \mu_D \circ T(h \circ_\D g) \circ f
      = \mu_D \circ T(\mu_D \circ Th \circ g) \circ f
      = \mu_D \circ T\mu_D \circ T^2h \circ Tg \circ f$$
      $$=^* \mu_D \circ \mu_{TD} \circ T^2h \circ Tg \circ f
      =^{**}\mu_D \circ Th \circ \mu_C \circ Tg \circ f
      =\mu_D \circ Th \circ (g \circ_\D f)
      = h \circ_\D (g \circ_\D f)$$

      Donde ($*$) es consecuencia del diagrama \eqref{diagrama:mu}
      de la definición de mónada aplicado sobre $D$ y ($**$) es
      consecuencia de la naturalidad de $\mu : \nat{T^2}{T}$.

    \item Existen identidades. Sean
      $f \in \Hom_\D(A, X)$ y $g \in \Hom_\D(X, A)$. Observamos
      que
      $$f \circ_\D \eta_A = \mu_X \circ Tf \circ \eta_A
      =^* \mu_X \circ \eta_{TB} \circ f
      =^{**} f$$
    Donde $(*)$ es consecuencia de la naturalidad de
    $\eta : \nat{1_\C}{T}$ y $(**)$ es consecuencia
    del diagrama \eqref{diagrama:eta} de la definición de mónada.
    Por otro lado
    $$\eta_A \circ_\D g = \mu_A \circ T\eta_A \circ g
    =^* g$$
    con $(*)$ consecuencia otra vez del diagrama \eqref{diagrama:eta}.

    Esto muestra que en la categoría $\D$ la flecha $\eta_A$ es la
    identidad del objeto $A$ para cualquier objeto $A$ de la categoría.
    \end{itemize}
  \item Ahora que conocemos la estructura de categoría
    de $\D$ definimos $F : \arr{\C}{\D}$ como $FA = A$
    para cada objeto $A$ de $\C$ y $Ff = Tf \circ \eta_A$ para
    toda flecha $f: \arr{A}{B}$ de $\C$.
    Veamos que $F$ es un funtor. Por partes:
    \begin{itemize}
    \item $F$ se lleva bien con la composición. Dado el diagrama
      $A \xrightarrow{f} B \xrightarrow{C}$ en la categoría $\C$
      tenemos que:
      $$F(g \circ f) = T(g \circ f) \circ \eta_A
      = Tg \circ Tf \circ \eta_A
      =^* Tg \circ \eta_B \circ f$$
      Donde $(*)$ se deduce de la naturalidad de
      $\eta : \nat{1_\C}{T}$. Igualmente:
      $$Fg \circ_\D Ff = (Tg \circ \eta_B) \circ_\D (Tf \circ \eta_A)
      = \mu_C \circ T(Tg \circ \eta_B) \circ (Tf \circ \eta_A)$$
      $$= \mu_C \circ T^2g \circ T\eta_B \circ Tf \circ \eta_A
      =^* \mu_C \circ T^2g \circ T\eta_B \circ \eta_B \circ f$$
      $$=^{**} Tg \circ \mu_B \circ T\eta_B \circ \eta_B \circ f
      =^{***} Tg \circ \eta_B \circ f$$
      Con $(*)$ consecuencia de la naturalidad de
      $\eta: \nat{1_\C}{T}$, $(**)$ consecuencia de la naturalidad
      de $\mu : \nat{T^2}{T}$, y $(***)$ por el diagrama
      \eqref{diagrama:eta} de la definición de mónadas.


    \item $F$ lleva identidades a identidades:
      $$F1_A = T1_A \circ \eta_A = 1_{TA} \circ \eta_A = \eta_A$$
      que es la flecha identidad de $A$ en $\D$.

    \end{itemize}

  \item Definimos $GA = TA$ para cada objeto $A$ de $\D$
    y $Gf = \mu_B \circ Tf$ para cada flecha $f \in \Hom_\D(A, B)$.
    Mostramos que $G: \arr{\D}{\C}$ es un funtor:
    \begin{itemize}
    \item $G$ se lleva bien con la composición. Dadas $f \in \Hom_\D(A, B)$
      y $g \in \Hom_\D(B, C)$ tenemos que:
      $$G(g \circ_\D f) = \mu_C \circ T(g \circ_\D f)
      = \mu_C \circ T(\mu_C \circ Tg \circ f)
      = \mu_C \circ T\mu_C \circ T^2g \circ Tf$$
      $$=^* \mu_C \circ \mu_{TC} \circ T^2g \circ Tf
      =^{**} \mu_C \circ Tg \circ \mu_B \circ Tf
      = Gg \circ Gf$$
      Con $(*)$ consecuencia del diagrama \eqref{diagrama:mu}
      de la definición de mónada y $(**)$ consecuencia de la
      naturalidad de $\mu : \nat{T^2}{T}$.
    \item $G$ se lleva bien con las identidades.
      $G\eta_A = \mu_A \circ T\eta_A = 1_{TA}$
      por el diagrama \eqref{diagrama:eta} de la definición de mónada.
    \end{itemize}
  \item La prueba de que $F \adjoint G$ es sencilla.
    El isomorfismo natural que necesitamos para establecer la
    adjunción es precisamente:
    $$\phi : \Hom_\D(F-, -) \cong \Hom_\C(-, G-)$$
    $$\phi_{A, B}(f) = f$$
    La comprobación de la naturalidad de $\phi$ es rutinaria.

  \end{enumerate}
\end{proof}


\section{En Programación}
\subsection{Adjunciones}
\paragraph{Currificación}
Definimos el siguiente funtor producto en Haskell.
\begin{verbatim}
type Product b a = (b, a)

instance Functor (Product fijo) where
  -- fmap :: a -> c -> (fijo, a) -> (fijo, c)
  fmap f (x, a) = (x, f a)
\end{verbatim}

Veamos que los funtores \cod{Product fijo} y y \cod{Reader fijo}
son funtores adjuntos. Para comprobarlo necesitamos definir
una biyección natural entre los conjuntos
$\Hom_\Hask(\cod{Product fijo a}, z) \cong
\Hom_\Hask(a, \cod{Reader fijo z})$. La biyección se puede
implementar de la siguiente forma:
\begin{verbatim}
biy :: ((Product fijo a) -> z) -> (a -> (Reader fijo z))
biy f a = Reader (\fijo -> f (a, fijo))

-- y la inversa
biy_inv :: (a -> Reader fijo z) -> (Product fijo a -> z)
biy_inv f (fijo, a) = (f a) fijo
\end{verbatim}
Se puede comprobar que son inversas fácilmente. La unidad la podemos
implementar tambień:
\begin{verbatim}
unit :: a -> Reader fijo (Product fijo a)
unit a = Reader (\fijo -> (fijo, a))
\end{verbatim}
Y la counidad:
\begin{verbatim}
counit :: Product fijo (Reader fijo a) -> a
counit (fijo, Reader f) = f fijo
\end{verbatim}

Y sabemos que son naturales porque son polimórficos en la variable
\cod{a}.

Esta adjunción nos dice una cosa sobre Haskell: es equivalente
dar una función que recibe dos parámetros de tipo
\cod{A} y de tipo \cod{B} y devuelve uno de tipo \cod{Z} que dar
una función que recibe un parámetro de tipo \cod{A} y devuelve
una función de tipo \cod{B -> Z}. Esto motiva la sintaxis de
Haskell para llamar a función. Habitualmente cuando en un lenguaje
se implementa una función \cod{f} que recibe dos parámetros ha
de ser llamada como:
\begin{verbatim}
f(a, b)
\end{verbatim}
En Haskell esto no es así. Debido a que \textit{es lo mismo} definir
una función que recibe dos parámetros que una función que recibe
un parámetro y devuelve otra función tenemos que Haskell opta por
la segunda. Entonces la sintaxis de Haskell

\begin{verbatim}
f a b = (f a) b
\end{verbatim}

Es decir, se aplica la función \cod{f} al valor \cod{a}, lo que devuelve
una función que recibe un parámetro \cod{b} y esta se aplica a continuación
al valor \cod{b}.
\textcolor{red}{Esto es útil en Haskell por varios motivos que explicitaremos}.

\subsection{Mónadas}
Al igual que existe una typeclass \cod{Functor} también existe una typeclass
\cod{Monad} en la bibiloteca estándar de Haskell. La definición es la siguiente:

\begin{verbatim}
class Monad M where
  return :: a -> M a
  bind :: M a -> (a -> M b) -> M b
\end{verbatim}

Esta definición de Mónada es similar a la que dimos originalmente. Podemos ver
que \cod{return} es precisamente la transformación natural $\eta$.
¿Pero cómo se relaciona \cod{bind} con la terminología que usamos anteriormente?
Implementamos la siguiente función:

\begin{verbatim}
mu :: Monad m => m (m a) -> m a
mu mma = bind mma id
\end{verbatim}

\subsubsection{Notación \cod{do}}


\chapter{Patrones de diseño}
En esta sección estudiaremos tres ejemplos de patrones de diseño
y conjuntos de
buenas prácticas que se encuentran habitualmente en código escrito
en haskell y que son inicialmente motivados por la teoría de
categorías.

\section{El patrón categórico}
La filosofía Unix, originada
principios de los setenta durante el desarrollo del sistema
operativo homónimo, aboga por el
diseño de componentes de software con una superficie muy reducida
y altamente combinable. Peter Salus resumió la filosofía Unix
en tres puntos:
\begin{itemize}
\item Escribe programas que hagan una sola cosa y que la hagan bien.
\item Escribe programas para que puedan funciona conjuntamente.
\item Escribe programas que trabajen con texto, porque el texto
  es un formato que entienden todos los programas.
\end{itemize}
La vigencia de la filosofía Unix aproximadamente 50 años más
tarde de su concepción es muestra de su potencia. La modularidad
es también una cualidad que se desea del software y para la que
se desarrollan patrones de diseño.

La programación funcional es otro ejemplo de ambiente en el que
se sigue esta filosofía, ya que se busca crear software
a base de componer de forma arbitrariamente elaborada
funciones sencillas.

El patrón categórico es un refinamiento de estas ideas, buscando
que las primitivas básicas que comprende tu software sean las
flechas de una categoría. Esto tiene las siguientes ventajas:
\begin{itemize}
\item Exista una forma de componer las primitivas. Tal como
  en la filosofía unix.
\item La composición es asociativa. Es mucho más sencillo razonar
  sobre una combinación de elementos asociativa que sobre una
  que no lo es. La asociatividad hace más sencillo comprender
  qué hace el código.

\item La existencia de identidades supone la existencia de primitivas
  que son \textit{triviales} respecto a la composición. Este es
  otro factor que favorece
  la comprensión de la operación de composición.
\end{itemize}

Estudiaremos un ejemplo de librería en haskell que fue concebida
con el patrón categórico en mente.

\subsection{Pipes}
\cod{pipes} es una librería de software libre hecha en haskell para
la gestión de \cod{streams}. Esta librería provee abstracciones para
el tratamiento de secuencias de datos y transformaciones sobre estas.
Un ejemplo de secuencias que podrían gestionarse con la librería
\cod{pipes} sería una sucesión de mensajes que llegan de un servidor
remoto o las líneas de un fichero en disco.

La librería ofrece tres abstracciones principales para el manejo
de \cod{streams}:
\begin{itemize}
\item \cod{Producers}: a los que en el presente texto nos referiremos
  como productores. Modelan fuentes de datos. Un \cod{Producer},
  por ejemplo, podría ser el que emite las líneas de un fichero de
  datos para su consumo por otra parte del software.
\item \cod{Consumers}: modelan destinos de datos. \textit{Piden} datos
  para realizar tareas con ellos. Un ejemplo de \cod{Consumer}
  recibiría mensajes de tipo \cod{String} y los almacenaría
  en una base de datos.
\item \cod{Pipes}: modelan transformaciones de los datos. Pueden recibir
  datos de un tipo y emitir datos de otro tipo. Un ejemplo sería un
  \cod{pipe} que recibe enteros y emite esos mismos enteros multiplicados
  por dos.
\end{itemize}

En este texto nos centraremos en los productores y pondremos de manifiesto
la naturaleza categórica en la que se centra su diseño. Explicaciones
análogas se pueden ofrecer para los consumidores y para los \cod{pipes}.

\subsubsection{Productores}
El tipo que describe a los productores es:
\begin{minted}{haskell}
Producer tipoEmitido monadaBase tipoResultado
\end{minted}
En lo que concierne a los ejemplos que mostraremos a lo largo de esta
sección consideraremos que la mónada base es \cod{IO}. Esto permite al
productor realizar acciones de entrada/salida como escribir en ficheros,
mandar información por la red o mostrar información por pantalla. El
\cod{tipoEmitido} es el tipo de los mensajes que emite el productor.
Cuando el productor acaba de emitir los mensajes que haya
considerado necesarios puede optar por devolver un valor. El tipo de
ese valor es \cod{tipoResultado}. El siguiente es un ejemplo de
productor con tipos concretos.
\begin{minted}{haskell}
threeNumbers :: Producer Int IO ()
threeNumbers = do
  yield 3
  yield 4
  lift $ putStrLn "Accion de IO"
  yield 5
  return ()
\end{minted}
La función \cod{lift} nos permite ejecutar acciones de la mónada
base en el código del productor. En este caso esa acción es mostrar
un mensaje por pantalla. Este productor emitirá tres valores enteros
y entre el segundo y el tercer valor ejecutará una acción de
entrada/salida.

Existe un tipo especial de productor que no emite ningún valor,
y se llama \cod{Effect}. En este texto nos referiremos a los valores
de tipo \cod{Effect} como efectos. \cod{pipes} no nos permite ejecutar
un productor si no lo \textit{conectamos} a algún sitio, es decir,
si no hay alguna parte de nuestro código que utilice los valores
emitidos. No existe este problema con los efectos, ya que no emiten
ningún valor. Para ejecutar un efecto podemos usar la función:
\begin{minted}{haskell}
runEffect :: Effect IO () -> IO ()
\end{minted}

Consideremos el efecto:
\begin{minted}{haskell}
print :: Int -> Effect IO ()
print x = do
  lift (putStrLn $ "el numero es " ++ show x)
\end{minted}
Esta función recibe como parámetro un entero y devuelve un efecto que
consiste en imprimir por pantalla tal entero. Podemos \textit{conectar}
esta función que recibe un entero con nuestro productor anterior, que
emitía enteros. La forma de combinarlos es usar la función \cod{for}
de \cod{pipes} que podemos considerar
que tiene la siguiente cabecera:
\begin{minted}{haskell}
for :: Monad m => Producer a m r -> (a -> Effect m ()) -> Effect m r
\end{minted}
Entonces podemos construir la expresión:
\begin{minted}{haskell}
(for threeNumbers print) :: Effect IO ()
\end{minted}
La forma en la que \cod{for} combinará el productor con la función
\cod{print} es sustituyendo cada llamada a \cod{yield} en el productor
por una llamada a la función \cod{print}. Es decir, la expresión anterior
generaría un efecto equivalente a:
\begin{minted}{haskell}
do
  lift (putStrLn $ "el numero es " ++ show 3)
  lift (putStrLn $ "el numero es " ++ show 4)
  lift $ putStrLn "Accion de IO"
  lift (putStrLn $ "el numeroes " ++ show 5)
  return ()
\end{minted}

En general, \cod{for} no solo permite combinar productores con
efectos sino que también se pueden combinar productores con otros
productores. Definimos la siguiente función:
\begin{minted}{haskell}
dup :: Int -> Producer Int IO ()
dup x = do
  yield x
  yield x
\end{minted}
es decir, recibe un parámetro y devuelve un productor que emite
ese valor dos veces. Este productor también se puede combinar con
\cod{for} siguiendo la misma lógica que se seguía anteriormente:
cada llamada a \cod{yield} se sustituye por una llamada a la función
\cod{dup}. Entonces el productor
\begin{minted}{haskell}
for threeNumbers dup
\end{minted}
es equivalente al productor
\begin{minted}{haskell}
do
  yield 3
  yield 3
  yield 4
  yield 4
  lift $ putStrLn "Accion de IO"
  yield 5
  yield 5
  return ()
\end{minted}
¿Podemos combinar también \cod{dup} y \cod{print}? La respuesta es
afirmativa. En \cod{pipes} se define un combinador con el siguiente
tipo:
\begin{minted}{haskell}
(~>) :: Monad m
     => (a -> Producer b m r)
     -> (b -> Producer c m r)
     -> (a -> Producer c m r)
(f ~> g) x = for (f x) g
\end{minted}
Es decir si tenemos una función que recibe un parámetro de tipo \cod{a}
y devuelve un productor que emite valores de tipo \cod{b} y tenemos
otra función que acepta un parámetro de tipo \cod{b} y emite valores
de tipo \cod{c} podemos combinarlos para obtener una función que
recibe un parámetro de tipo \cod{a} y devuelve un productor que emite
valores de tipo \cod{c}. Esta cabecera puede resultar confusa para
alguien que no haya programado nunca en haskell por lo que proponemos
el siguiente ejemplo:
\begin{minted}{haskell}
dup2 :: String -> Producer String IO ()
dup2 s = do
  yield s
  yield (s ++ s)

leng :: String -> Producer Int IO ()
leng s = yield (length s)

numberAndSquare :: Int -> Producer Int IO ()
numberAndSquare x = do
  yield x
  lift $ putStrLn "Ahora emitimos el cuadrado"
  yield (x * x)
\end{minted}
El operador \cod{\pcomp >} nos permite combinar tanto \cod{dup2} y
\cod{leng} como \cod{leng} y \cod{numberAndSquare}. Detallamos
qué significa cada una de las composiciones:
\begin{minted}{haskell}
dup2 ~> leng :: String -> Producer Int IO ()
\end{minted}
es una función equivalente a:
\begin{minted}{haskell}
dup2Leng :: String -> Producer Int IO ()
dup2Leng  = do
  yield (length s)
  yield (length (s ++ s))
\end{minted}
Por otro lado:
\begin{minted}{haskell}
leng ~> numberAndSquare :: String -> Producer Int IO ()
\end{minted}
es equivalente a:
\begin{minted}{haskell}
lengNumberAndSquare :: String -> Producer Int IO ()
lengNumberAndSquare s = do
  yield (length s)
  lift $ putStrLn "Ahora emitimos el cuadrado"
  yield ((length s) * (length s))
\end{minted}
Pero aún hay más, también son posibles las siguientes combinaciones:
\begin{minted}{haskell}
(dup2 ~> leng) ~> numberAndSquare
-- y
dup2 ~> (leng ~> numberAndSquare)
\end{minted}
y no solo eso sino que además a partir de la implementación de
\cod{pipes} se puede demostrar que las dos expresiones anteriores
producen productores equivalentes. En ese sentido
podemos decir que el operador
\cod{\pcomp >} es asociativo.

De hecho podemos decir que los productores forman una categoría.
En esta categoría los objetos son los tipos de haskell y las flechas
$\arr{a}{b}$ son las funciones de haskell de tipo
\cod{a -> Producer b IO ()}. La composición es el operador
\cod{\pcomp >} que hemos descrito anteriormente. Ya hemos visto que el operador
es asociativo, quedaría ver que existen las identidades para cada
tipo y existen. La identidad es precisamente
\begin{minted}{haskell}
yield :: a -> Producer a IO ()
\end{minted}
es decir, \cod{yield \pcomp> f = f} y \cod{g \pcomp> yield = g}.

\section{Monad Transformers}
En el capítulo anterior comentamos el uso de mónadas para realizar
labores como:
\begin{itemize}
\item Ejecutar operaciones de entrada/salida
\item Gestionar computaciones puras con estado
\item Gestión de fallos de forma transparente
\item Acceso global a un repositorio de información inmutable
\end{itemize}
Una aplicación normalmente necesita varias de estas funcionalidades
al mismo tiempo. Los \textit{monad transformers} permiten combinar
las funcionalidades de varias mónadas. Un ejemplo de clase que
regula el funcionamiento de esta abstracción es:

\begin{minted}{haskell}
class MonadTrans t where
    lift :: Monad m => m a -> t m a
\end{minted}
Es decir, un \cod{t} t

\subsection{Ejemplo de aplicación}
Podemos suponer que estamos escribiendo una aplicación que necesita
acceso global a unos datos de configuración, tiene que realizar
transformaciones que modifiquen estado y además tiene que mostrar
información por pantalla. Si tuviéramos que realizar esas tareas
por separado usaríamos respectivamente el \cod{Reader} monad,
el \cod{State} monad y el \cod{IO} monad. Para hacer todo a la vez
tendremos que usar monad transformers.

Nuestra aplicación hará lo siguiente: recibirá por entrada estándar
enteros que escriba el usuario y los irá almacenando en una lista.
Cuando se acabe la entrada mostrará la lista completa de los enteros
que el usuario ha introducido. Es una aplicación sencilla pero que
mostrará la utilidad de usar monad transformers.

En primer lugar definiremos el tipo \cod{Config} de la aplicación que
representa la configuración inicial de la aplicación. En este caso
el tipo \cod{Config} solo almacenará el título de la aplicación y
el carácter prompt que usará a la hora de mostrar salida. El tipo
es el siguiente:
\begin{minted}{haskell}
data Config = Config {
  title :: String,
  prompt :: String
  }
\end{minted}
El estado a modificar de nuestra aplicación será precisamente la
lista de enteros. Usaremos un alias de tipo para representar este
tipo.
\begin{minted}{haskell}
type Stack = [Int]
\end{minted}
Definimos la pila de mónadas sobre la que diseñaremos nuestra
aplicación. Usaremos una mónada
\cod{Reader Config} para que las acciones
que se ejecuten en el dominio de nuestra aplicación
tengan acceso a la configuración global. Utilizaremos también
una mónada \cod{State Stack} para que las acciones puedan modificar
el estado de la lista de enteros. Finalmente, utilizaremos la
mónada \cod{IO} para poder efectuar acciones de entrada/salida.
El tipo que utilizaremos para definir las acciones de nuestra aplicación
será entonces:
\begin{minted}{haskell}
type AppT = ReaderT Config (StateT Stack IO)

newtype AppM a = AppM { unApp :: AppT a }
  deriving (Functor, Applicative, Monad,
            MonadReader Config, MonadState Stack, MonadIO)
\end{minted}
El tipo \cod{AppT} es un alias para la pila de monad transformers
que usaremos. Encapsulamos ese tipo en otro llamado \cod{AppM} (esto
es algo habitual en haskell y representa ventajas a la
hora de aumentar la seguridad en cuanto a la comprobación de tipos
que el compilador puede hacer). La librería \cod{mtl} nos proporciona
typeclasses como \cod{MonadReader Config} que nos permiten utilizar
la funcionalidad del monad \cod{Reader} dentro de una pila
de monad transformers. Análogamente sucede con
\cod{MonadState Stack}.

Si tenemos un valor de \cod{AppM a} necesitamos una forma de
transformarlo en un valor de \cod{IO a} para poder ejecutarlo.
Para ello usamos la función:
\begin{minted}{haskell}
runApp :: Config -> AppM a -> IO a
runApp config (AppM m) = evalStateT (runReaderT m config) []
\end{minted}
Esta aplicación recibe una configuración y un valor de
\cod{AppM a} y devuelve la acción de \cod{IO a} asociada. Ejecuta
la acción inicializando el valor del \cod{Stack} a la lista vacía.

Mostramos algunas de las acciones que se ejecutarán en el contexto
de nuestra aplicación:
\begin{minted}{haskell}
addNumber ::
  (MonadState Stack m) => Int -> m ()
addNumber x = modify $ \state -> (x:state)
\end{minted}
Esta acción añade al estado de la aplicación un entero \cod{x}. Nótese
que esta acción se podría implementar con este tipo y sería equivalente:
\begin{minted}{haskell}
addNumber ::
  Int -> AppM ()
addNumber x = modify $ \state -> (x:state)
\end{minted}
La ventaja de implementarla con un tipo polimórfico que sea
\cod{MonadState Stack m} es que el compilador entiende que esa función
puede ejecutarse sobre cualquier pila de mónadas que gestione un estado
de tipo \cod{Stack}. Esto no es útil porque existan muchas pilas de
mónadas sobre la que sea útil ejecutar esa acción sino porque con
ese tipo se hace imposible ejecutar ninguna acción \cod{IO ()} en
la función. Para que se pueda ejecutar una acción de entrada/salida
se debe estar dentro de un \cod{MonadIO} pero la \textit{constraint}
\cod{MonadState Stack m} no implica que se pueda ejecutar ninguna
operación \cod{IO}. La ventaja de este hecho es que viendo
simplemente el tipo
de la función podemos saber que la única funcionalidad que puede
ejecutar esta función es una transformación sobre el estado
\cod{Stack}. En este caso es fácil leer directamente qué hace la función
y saber que solo modifica el estado, pero delimitar estrictamente
sobre qué partes de un sistema actúa una función es útil a la hora
de trabajar sobre sistemas más complejos.

Otro ejemplo de esto lo encontramos en la acción:
\begin{minted}{haskell}
appPrint ::
  (MonadReader Config m, MonadIO m) => String -> m ()
appPrint msg = do
  promptText <- fmap prompt ask
  liftIO $ putStrLn (promptText ++ " " ++ msg)
\end{minted}
Viendo este código sabemos que la función \cod{appPrint} no puede
modificar el estado. En las constraints asumimos \cod{MonadIO}
porque necesitamos imprimir por la salida estándar y asumimos
\cod{MonadReader} porque necesitamos leer la configuración
para ver el carácter de \cod{prompt}.

\begin{minted}{haskell}
showTitle ::
  (MonadReader Config m, MonadIO m) => m ()
showTitle = do
  titleStr <- fmap title ask
  appPrint titleStr
\end{minted}
Es análoga a \cod{appPrint}. Implementamos también una función
que muestra por pantalla el estado de la aplicación.
\begin{minted}{haskell}
showStack :: AppM ()
showStack = do
  stack <- get
  appPrint (show stack)
\end{minted}
En esta función no usamos un tipo polimórfico porque necesitamos
utilizar las tres mónadas para cumplir esta funcionalidad. Necesitamos
el \cod{State} monad para leer el estado de la aplicación, necesitamos
\cod{Reader} para poder usar \cod{appPrint}, que necesita conocer
el \cod{prompt} y por supuesto necesitamos \cod{MonadIO} para imprimir
mensajes por pantalla.

Finalmente mostramos el código que falta para que la aplicación
funcione:
\begin{minted}{haskell}
appLoop :: AppM ()
appLoop = do
  eof <- liftIO isEOF
  if eof
    then showStack
    else (do n <- liftIO readLn
             addNumber n
             appLoop
         )

main :: IO ()
main = runApp (Config "AppTitle" ">") $ do
  showTitle
  appLoop
\end{minted}
La función \cod{main} inicializa la aplicación con un valor
de configuración, imrpime el título por pantalla y ejecuta el
bucle central de la aplicación. Este bucle comprueba si se
ha llegado al final de la entrada estándar. Si no, se lee un entero
por la entrada, se añade a la pila y se vuelve a ejecutar el bucle
central. Si se ha llegado al final de la entrada estándar se
muestra por pantalla el contenido de la pila y se sale.

\section{Otros patrones de diseño}
Finalizaremos el trabajo con la exposición breve de un par de
patrones de diseño conocidos y bien documentados en la comunidad
de haskell que muestran que la influencia de la teoría de categorías
no se limita al patrón categórico y a los monads transformers.
Analizar en más profundidad las siguientes construcciones podría
ser una buena forma de continuar con la línea de este trabajo.

\subsection{\cod{Applicative}}
\cod{Applicative} es otra typeclass de la biblioteca estándar
de haskell. Es una superclase de \cod{Monad}, en el sentido de que
toda instancia de \cod{Monad} es una instancia de
\cod{Applicative}. La interfaz de \cod{Applicative} es la siguiente:
\begin{minted}{haskell}
class Functor f => Applicative f where
    pure :: a -> f a
    (<*>) :: f (a -> b) -> f a -> f b
\end{minted}
La función \cod{pure} nos recuerda a la función \cod{return} de la
instancia de \cod{Monad}, y de hecho si un applicative dado es
un \cod{Monad} las funciones \cod{pure} y \cod{return} deben coincidir.
La función \cod{(<*>)} se parece tanto a \cod{fmap} de la clase
\cod{Functor} como a \cod{(>{}>=)} de la clase \cod{Monad}, pero no
es ninguna de las dos. Veamos las leyes que una instancia de
\cod{Applicative} debe cumplir según la documentación de haskell:
\begin{minted}{haskell}
pure id <*> v = v -- identidad
pure (.) <*> u <*> v <*> w = v <*> (v <*> w) -- composición
pure f <*> pure x = pure (f x) -- homomorfismo
u <*> pure y = pure (\f -> f y) <*> u -- intercambio
\end{minted}
Se puede probar que si tenemos una instancia de mónada y
definimos:
\begin{minted}{haskell}
pure = return

-- (<*>) :: f (a -> b) -> f a -> f b
ff <*> fa = do
  f <- ff
  a <- fa
  return (f a)
\end{minted}
(hemos podido usar notación \cod{do} porque estamos asumiendo
que \cod{f} es un \cod{Monad})
Entonces se cumplen las leyes de applicative. A partir de
una instancia de applicative también podemos definir una instancia
de funtor que cumple las leyes de la siguiente forma:
\begin{minted}{haskell}
fmap :: (a -> b) -> f a -> f b
fmap f fa = (pure f) <*> fa
\end{minted}

Para entender la diferencia cualitativa entre \cod{Applicative} y
\cod{Monad} enunciaremos el siguiente resultado: toda expresión
creada usando los combinadores aplicativos (\cod{pure} y \cod{(<*>)})
se puede transformar en una expresión de la siguiente forma:
\begin{minted}{haskell}
pure f <*> u1 <*> u2 <*> ... <*> un
\end{minted}
donde \cod{f :: u1 -> u2 -> ... -> z} y \cod{ui :: f ui} son valores
aplicativos. Esto no es posible en general con expresiones monádicas,
donde la \textit{estructura} de la computación puede depender de
resultados intermedios.

Esta diferencia ha permitido que \cod{Applicative} se identifique
como la abstracción adecuada
para la resolución de múltiples problemas como el parseo de gramáticas libres de
contexto o la validación de formularios. Facebook liberó
recientemente
una librería de haskell llamada Haxl para la gestión
de la concurrencia de
forma totalmente implícita para el programador, permitiendo
que funciones como:
\begin{minted}{haskell}
numCommonFriends :: User -> User -> Haxl Int
numCommonFriends x y =
    length <$> (intersect <$> friendsOf x <*> friendsOf y)
\end{minted}
ejecuten las llamadas \cod{friendsOf} a la base de datos de forma
concurrente (e incluso integrando soluciones de cacheo de forma
transparente para el programador). Usando la instancia de
applicative de Haxl el programador permite que la librería
extraiga el paralelismo implícito de su código sin tener que programar
de forma distinta a como lo haría con código secuencial.

Categóricamente, \cod{applicative} ha sido identificado como una
clase especial de funtores de categorías monoidales.

\subsubsection{Free Monads}
Los free monads son una abstracción que proviene también de la
teoría de categorías y que ha sido utilizada exitosamente para el
diseño de software. La aplicación de las free monads a haskell
permite el encaje natural de lenguajes de dominio específico así
como la creación de intérpretes muy generales para estos lenguajes.

Una de las librerías más utilizada que implementan la construcción
es \cod{free}, y es utilizada por numerosas aplicaciones en haskell
entre las que podemos destacar:
\begin{itemize}
\item \cod{engine-io}: una implementación de un protocolo de
  websockets.
\item \cod{reddit}: una librería para interactuar con el popular
  servicio homónimo.
\end{itemize}

La idea que sugiere la utilización de las free monads es la
creación de una DSL que permita expresar toda la lógica de negocio
de tu aplicación de forma agnóstica al resto de tecnologías
utilizadas (es decir, no mezclar código del framework que estés
utilizando con tu lógica de negocio) y que luego sean los intérpretes
de esa DSL los que tengan que interactuar con el resto de componentes
del sistema.


\chapter*{Conclusiones y vías futuras}
En la asignatura de Geometría I, durante mi primer año de carrera
pude leer una afirmación que despertó mi curiosidad: ``
\textit{un espacio vectorial es naturalmente isomorfo a su doble dual}
''. Lo que
entendí en su momento es que entre esos dos espacios existía
un isomorfismo que era más evidente que otros, aunque no creí
que esa noción se pudiera formalizar.

Este trabajo, cinco años más tarde, me permite por fin
darle un sentido preciso a esa afirmación. Digo más, la maquinaria
que me permite entender qué significa que un espacio vectorial
sea naturalmente isomorfo a su doble dual me
sugiere una forma nueva de entender gran parte de los contenidos
de la carrera, especialmente (pero no exclusivamente) las asignaturas
de álgebra.

La complejidad del software construido a lo larga de la escasa
historia de la disciplina no ha hecho más que crecer, y no hay motivos
para pensar que esta tendencia parará. El desarrollo de
nuevos paradigmas como la programación
orientada a objetos permitió a los ingenieros domar la complejidad
y llevar a cabo proyectos que hubieran sido imposibles con técnicas
más rudimentarias. La utilización de la teoría de categorías
para formalizar abstracciones en lenguajes de programación funcional
podría ser un paso adelante. Muchas empresas se benefician ya
del uso de haskell en sus sistemas y de las
numerosas innovaciones que este
lenguaje de programación introdujo.

Desde mi punto de vista he hecho una buena exposición de los contenidos
de teoría de categorías que permiten entender los patrones
de diseño más usados en haskell. No estoy convencido de que este
texto, por sí mismo, pueda convencer a nadie de que estas técnicas
pueden ser beneficiosas para el desarrollo de software. Haskell es
un lenguaje de programación muy diferente, y cómo tal, aceptar sus
propuestas puede requerir un trabajo continuado. En este sentido
podemos afirmar que los objetivos han sido alcanzados.

La línea exacta de este trabajo puede continuarse con el estudio
de mónadas libres y categorías de álgebras sobre mónadas. Estos
conceptos, que se estudian en los textos de categorías por su
relevancia en las matemáticas, tienen aplicaciones
a la programación. Otra línea de trabajo sería la de la teoría
de categorías para el estudio de sistemas de tipos.

\printbibliography

\appendix
\chapter[Haskell Básico]{Haskell básico para entender los ejemplos}
Comentamos brevemente algunas funcionalidades de haskell
que facilitarán la comprensión de los ejemplos utilizados.

\section{Expresiones y tipos}
Toda expresión en Haskell tiene un tipo. Los siguientes
son ejemplos de expresiones junto con sus tipos:
\begin{minted}{haskell}
True :: Bool
False :: Bool
'c' :: Char
"hola" :: String
3 :: Int
\end{minted}
En general la sintaxis \cod{a :: A} se utiliza en haskell para
afirmar que la expresión \cod{a} tiene tipo \cod{A}. Que toda
expresión en haskell esté tipada no significa que tengamos que
ser explícito acerca de estos tipos, el compilador es capaz
en la mayoría de las ocasiones de inferir por nosotros cual
es el tipo de todas las expresiones de un programa. Será habitual,
sin embargo, que especifiquemos un tipo aunque no sea necesario,
para ayudar al lector.

\section{Funciones}
En haskell las funciones son ciudadanos de primera clase
(\textit{first class citizens}) en el sentido de que son valores
como cualquier otro. Esto significa, entre otras cosas, que las
funciones tienen un tipo.

Proponemos ejemplos de funciones a continuación:
\begin{minted}{haskell}
suma1 :: Int -> Int
suma1 x = x + 1

tres :: Int
tres = suma 2

esA :: Char -> Bool
esA c = if c == 'A'
        then True
        else False
\end{minted}
El tipo de \cod{suma1} indica que la función recibe un parámetro
de tipo entero y devuelve un valor de tipo entero. Se pueden
definir también funciones de múltiples parámetros:
\begin{minted}{haskell}
suma :: Int -> Int -> Int
suma x y = x + y

cinco :: Int
cinco = suma 3 2
\end{minted}
El tipo de \cod{suma} puede resulta sorprendente
a quien esté acostumbrado a trabajar con lenguajes imperativos.
Haskell solo considera funciones de un parámetro. La trampa es que
\cod{suma} es una función que recibe un parámetro de tipo \cod{Int}
y devuelve \emph{otra función} que recibe un parámetro \cod{Int}
y devuelve un valor de tipo \cod{Int}. El siguiente ejemplo es
equivalente al anterior:
\begin{minted}{haskell}
suma :: Int -> Int -> Int
suma x y = suma x + y

sumaTres :: Int -> Int
sumaTres = suma 3

cinco :: Int
cinco = sumaTres 2
\end{minted}
Entonces cuando haskell evalúa la expresión
\begin{minted}{haskell}
suma 3 2
\end{minted}
lo está haciendo asociado a la izquierda:
\begin{minted}{haskell}
(suma 3) 2
\end{minted}
La técnica de transformar una función que recibe varios parámetros
en una función que recibe un parámetro y devuelve otra función
se llama currificación. Haskell hace currificación automática
de todas las funciones.

De hecho no existe ninguna diferencia entre el tipo de la función
\cod{suma} tal como estaba definido anteriormente y el tipo:
\begin{minted}{haskell}
suma :: Int -> (Int -> Int)
suma x y = x + y
\end{minted}

\section{Expresiones lambda}
Existe un tipo especial de expresión en haskell llamado expresión
lambda y que sirve para definir funciones. Ponemos un ejemplo
a continuación:
\begin{minted}{haskell}
suma1 :: Int -> Int
suma1 = (\x -> x + 1)
\end{minted}
La expresión construye la función que recibe un valor \cod{x}
y devuelve el valor \cod{x+1}. Las expresiones lambda sirven
también para definir funciones que reciben múltiples parámetros:
\begin{minted}{haskell}
multiplica :: Int -> Int -> Int
multiplica = (\x y -> x * y)
\end{minted}

Es importante notar que la expresión \cod{\\x -> f x} es
equivalente a la expresión \cod{f}.

\section{Definir tipos}
Haskell permite al usuario definir sus propios tipos. Veamos
un ejemplo sencillo:
\begin{minted}{haskell}
data Dni = Dni Int Char
\end{minted}
Esta línea define el tipo \cod{Dni} (izquierda) y el constructor
homónimo (derecha) que acepta como parámetros un entero y un carácter.
Mostramos a continuación ejemplos de valores del tipo \cod{Dni}.
\begin{minted}{haskell}
(Dni 44669279 'J') :: Dni
(Dni 12345678 'C') :: Dni
\end{minted}
Todos los valores del tipo \cod{Dni} son de la forma
\cod{(Dni i l)} donde \cod{i :: Int} y \cod{l :: Char}. Además,
podemos considerar al constructor \cod{Dni} una función con tipo
\cod{Dni :: Int -> Char -> Dni}. Nótese que
\cod{Dni} antes de los dos puntos representa al constructor pero
que el \cod{Dni} de la derecha representa al tipo.

Como todos los valores del tipo \cod{Dni} son de la forma especificada
arriba, haskell nos permite definir funciones sobre ese tipo de la
siguiente manera:
\begin{minted}{haskell}
suma3AlNumero :: Dni -> Int
suma3AlNumero (Dni numero caracter) = numero + 3

-- devuelve el dni completo como una cadena
devuelveCadena :: Dni -> String
devuelveCadena (Dni n c) = (show n) ++ (show c)
\end{minted}

Podemos definir tipos con más de un constructor. Por ejemplo:
\begin{minted}{haskell}
data PersonaConAltura = Hombre Int | Mujer Int | Otro String Int
\end{minted}
donde definimos el tipo \cod{PersonaConAltura} y los constructores
\cod{Hombre :: Int -> PersonaConAltura},
\cod{Mujer :: Int -> PersonaConAltura},
\cod{Otro :: String -> Int -> PersonaConAltura}. Todo valor de
\cod{PersonaConAltura} será o bien de la forma
\cod{Hombre n}, o bien \cod{Mujer n}, o bien \cod{Otro o n}, donde
\cod{n :: Int} y \cod{o :: String}. Esto hace que haskell
nos permita definir funciones sobre este tipo como:
\begin{minted}{haskell}
comentarioSobreAltura :: PersonaConAltura -> String
comentarioSobreAltura (Hombre altura) = if altura < 174
                                        then "Tienes altura inferior a la media"
                                        else "Tienes altura superior a la media"
comentarioSobreAltura (Mujer altura) = if altura < 163
                                       then "Tienes altura inferior a la media"
                                       else "Tienes altura superior a la media"
comentarioSobreAltura (Otro genero x) =
            "No tengo informacion sobre la altura media en el género " ++ genero
\end{minted}
Esta funcionalidad de hacer distinción de casos sobre los constructores
de un tipo es conocida como \textit{pattern matching}. También se puede
hacer \textit{pattern matching} dentro de una expresión con la siguiente
sintaxis:
\begin{minted}{haskell}
genero :: PersonaConAltura -> String
genero persona = case persona of
    Hombre _ -> "Hombre"
    Mujer _ -> "Mujer"
    Otro g _ -> g

-- equivalentemente podríamos haber implementado la función como
genero (Hombre _) = "Hombre"
genero (Mujer _) = "Mujer"
genero (Otro g _) = g
\end{minted}
donde usamos el carácter \cod{\_} para indicar (tanto a lectores como
al compilador) que no usaremos esa variable.

Nótese que podemos definir tipos \textit{recursivos} en el siguiente
sentido:
\begin{minted}{haskell}
data Peano = Cero | Succ Peano
\end{minted}
Aquí estamos definiendo el tipo \cod{Peano} que tiene dos constructores:
\cod{Cero :: Peano} (no admite parámetros) y \cod{Succ :: Peano -> Peano}.
Podemos hacer \textit{pattern matching} sobre este tipo:
\begin{minted}{haskell}
aNumero :: Peano -> Int
aNumero Cero = 0
aNumero (Succ peano) = 1 + (aNumero peano)

esCero :: Peano -> Bool
esCero Cero = True
esCero (Succ _) = False
\end{minted}

Existe otra sintaxis para definir tipos con un solo constructor y que reciben
un solo parámetro:
\begin{minted}{haskell}
newtype EnvuelveEntero = EnvuelveEntero Int
\end{minted}
que para los propósitos de este trabajo consideraremos igual a la sintaxis
que usa \cod{data} en lugar de \cod{newtype}.

Cuando usamos la sintaxis:
\begin{minted}{haskell}
data Dni = Dni {
  numero :: Int
  letra :: Char
  }
\end{minted}
No solo estamos definiendo el tipo \cod{Dni} junto con el constructor de tipo
\cod{Dni :: Int -> Char -> Dni} sino que además estamos definiendo las
proyecciones evidente \cod{numero :: Dni -> Int} y \cod{letra :: Dni -> Char}.

\section{Funciones con tipos variables}
Podemos definir funciones polimórficas como la siguiente:
\begin{minted}{haskell}
aplicar :: (a -> b) -> a -> b
aplicar f x = f x

seis :: Int
seis = aplicar suma1 5

hombre :: String
hombre = aplicar genero (Hombre 180)
\end{minted}
La función \cod{aplicar} se puede especializar a tipos concretos como
\cod{a = Int, b = Int} en el caso del valor \cod{seis} o a los
tipos \cod{a = PersonaConAltura, b = String} en el caso del
valor \cod{hombre}.


\section{Tipos Parametrizados}
Podemos definir tipos parametrizados por otros tipos. Por ejemplo:
\begin{minted}{haskell}
data ParConNombre a = ParConNombre String a a
\end{minted}
esto define todos los tipos \cod{ParConNombre A} donde \cod{A} es
un tipo fijo pero arbitrario. Nótese que \cod{ParConNombre} no
es un tipo, pero \cod{ParConNombre Int} sí que lo es.
Un ejemplo de valor de este tipo sería \cod{ParConNombre "tresdos" 3 2 :: ParConNombre Int}.

El constructor \cod{ParConNombre :: String -> a -> a -> ParConNombre a}
es ahora una función con tipos variables. Se puede estar parametrizado por más
de un tipo:
\begin{minted}{haskell}
data Par a b = Par a b

-- ejemplos de valores
Par 3 'c' :: Par Int Char
Par 'c' 3 :: Par Char Int
Par 4 5 :: Par Int Int
\end{minted}

\section{Listas}
Podemos definir listas en haskell de forma similar a como
definimos el tipo \cod{Peano}:
\begin{minted}{haskell}
data Lista a = ListaVacia | NoVacia a Lista

-- ejemplos de valores
ListaVacia :: Lista Int
NoVacia 3 ListaVacia :: Lista Int
NoVacia 3 (NoVacia 4 ListaVacia) :: Lista Int
\end{minted}
y podemos definir funciones sobre este tipo:
\begin{minted}{haskell}
sumarLista :: Lista Int -> Int
sumarLista ListaVacia = 0
sumarLista (NoVacia x otraLista) = x + (sumarLista otraLista)
\end{minted}
No necesitamos definir este tipo por nuestra cuenta
porque ya está definido en la biblioteca estándar de haskell
y cuenta con una sintaxis más intuitiva:
\begin{minted}{haskell}
sumarLista :: [Int] -> Int
sumarLista [] = 0
sumarLista (x:xs) = x + (sumarLista xs)
\end{minted}
Cuando queramos hacer referencia a que \cod{[]} es un tipo
parametrizado a veces escribiremos \cod{[] Int} en lugar
de \cod{[Int]}.

\section{Typeclasses}
Haskell soporta una funcionalidad similar a las interfaces de java
conocida como typeclasses. Una typeclass define una interfaz que
cualquier tipo que quiera ser instancia de esa clase debe implementar.
Por ejemplo:
\begin{minted}{haskell}
data Tamaño = Grande | Pequeño

class TieneTamaño t where
  tamaño :: t -> Tamaño

instance TieneTamaño Int where
  tamaño x = if x > 179 then Grande else Pequeño

instance TieneTamaño String where
  tamaño s = if (length s) > 10 then Grande else Pequeño

instance TieneTamaño PersonaConAltura
  tamaño (Hombre altura) = tamaño altura
  tamaño (Mujer altura) = tamaño altura
  tamaño (Otro _ altura) = tamaño altura
\end{minted}
Nótese como hemos implementado la instancia de \cod{Tamaño} de
\cod{PersonaConAltura} aprovechando la instancia de \cod{Int}.
Podemos escribir funciones que solo acepten instancias de
esta typeclass:
\begin{minted}{haskell}
nombreTamaño :: Tamaño -> String
nombreTamaño Grande = "grande"
nombreTamaño Pequeño = "pequeño"

decirTamaño :: TieneTamaño a => String -> a -> String
decirTamaño nombre a =
  nombre ++ " tiene tamaño " ++ nombreTamaño (tamaño a)

-- ejemplos
decirTamaño "tres" 3 -- tres tiene tamaño pequeño
decirTamaño "Braulio" (Hombre 180) -- Braulio tiene tamaño grande
decirTamaño "cadena de texto" "eeeeeeeeeeee" -- ...
\end{minted}
De hecho, también se pueden definir clases que solo puedan ser
instanciadas por instancias de otras clases:
\begin{minted}{haskell}
class TieneTamaño t => TieneTamañoNumerico t where
  tamañoNumerico :: t -> Int

instance TieneTamañoNumerico Int where
  tamañoNumerico x = x

instance TieneTamañoNumerico Char where
  tamañoNumerico _ = 100 -- Error: Char no es instancia de TieneTamaño
\end{minted}

\end{document}
